\documentclass[11pt]{amsart}

\usepackage[T1]{fontenc}
\usepackage{lmodern}
\usepackage{amsmath,amssymb,mathtools}
\usepackage{microtype}
\usepackage[colorlinks=true,linkcolor=blue,citecolor=blue,urlcolor=blue,filecolor=blue]{hyperref}
\usepackage[nameinlink,noabbrev]{cleveref}

\newtheorem{theorem}{Theorem}[section]
\newtheorem{proposition}[theorem]{Proposition}
\newtheorem{lemma}[theorem]{Lemma}
\newtheorem{corollary}[theorem]{Corollary}
\theoremstyle{remark}
\newtheorem{remark}[theorem]{Remark}

\numberwithin{equation}{section}

\title[Subregular affine cells for $L_{-1}(D_\ell)$]{Subregular affine cells and the level $-1$ vertex algebra of type $D$}
\author{Sihai Jin}
\address{Department of Mathematics, Sichuan University, China}
\email{jinsihai@stu.scu.edu.cn}
\keywords{affine vertex algebras, affine Weyl groups, Kazhdan--Lusztig cells, Zhu algebras, finite $W$-algebras, primitive ideals}
\date{}

\begin{document}

\begin{abstract}
We prove the Shan--Yan--Zhao affine left-cell conjecture for the simple affine vertex algebras $L_{-1}(D_\ell)$, $\ell\ge5$.  The distinguished vacuum block has exactly $\ell+1$ simple objects, indexed by the subregular affine left cell containing $s_0$.  Two finite-dimensional ingredients enter the proof.  A primitive-ideal inclusion shows that the predicted non-vacuum modules descend to the relevant quotient, while a separate exhaustion argument rules out further highest weights.  It follows that the Grothendieck group is the $q=1$ specialization of the corresponding dual affine left-cell module.  The same identification yields uniform character formulas in terms of subregular inverse Kazhdan--Lusztig coefficients.
\end{abstract}

\maketitle

\section{Introduction and main results}

Let \(L_k(\mathfrak g)\) denote the simple affine vertex algebra associated
with a simple complex Lie algebra \(\mathfrak g\).  Shan--Yan--Zhao predict
that certain distinguished blocks of \(L_k(\mathfrak g)\)-modules are
controlled by affine Kazhdan--Lusztig left cells; in particular, the
Grothendieck group of such a block should be the \(q=1\) specialization of
the corresponding dual affine left-cell module \cite{ShanYanZhao2026}.  The
case treated here is
\[
   (\mathfrak g,k)=(D_\ell,-1),\qquad \ell\ge5.
\]

We use the Bourbaki realization
\[
 \alpha_i=\varepsilon_i-\varepsilon_{i+1}\ (1\le i\le\ell-1),
 \qquad \alpha_\ell=\varepsilon_{\ell-1}+\varepsilon_\ell,
\]
with \(\theta=\varepsilon_1+\varepsilon_2\) and
\(\rho=(\ell-1,\ell-2,\ldots,1,0)\).  Put
\[
   N:=k+h^\vee=2\ell-3,\qquad r:=\ell-2,
\]
and let \(s_0,s_1,\ldots,s_\ell\) be the simple reflections of
\(D_\ell^{(1)}\).  For \(1\le i\le\ell\), let \(z_i\) be the product of
finite simple reflections along the unique Dynkin path from \(i\) to \(2\),
ending in \(s_2\), and set
\[
   w_0=s_0,\qquad w_i=z_i s_0,\qquad
   \mu_0=0,\qquad \mu_i=z_i\rho-\rho.
\]
The subregular-cell description of
Bezrukavnikov--Kac--Krylov \cite{BezrukavnikovKacKrylov2024} gives
\[
   \mathbf c^L(s_0)=\{w_0,w_1,\ldots,w_\ell\},
   \qquad
   w_i\circ(-\Lambda_0)=-\Lambda_0+\mu_i.
\]

\begin{theorem}[Main theorem]
For every \(\ell\ge5\),
\begin{equation}
   \operatorname{Irr}\mathcal O_{-\Lambda_0}
   \!\left(L_{-1}(D_\ell)\right)
   =
   \left\{
      L\!\left(w_i\circ(-\Lambda_0)\right)
      \;\middle|\;0\le i\le\ell
   \right\}.
\end{equation}
Moreover, under the ambient realization of Shan--Yan--Zhao,
\begin{equation}
   K_0\mathcal O_{-\Lambda_0}\!\left(L_{-1}(D_\ell)\right)
   \cong
   \left.
      \mathcal H^\vee_{\mathrm{aff},\mathbf c^L(s_0)}
   \right|_{q=1}.
\end{equation}
Thus \cite[Conjecture~5.2.1]{ShanYanZhao2026} holds for
\((D_\ell,-1)\).  The \(\ell+1\) simple characters are given by the uniform
subregular inverse Kazhdan--Lusztig formula of
\cite{BezrukavnikovKacKrylov2024}; see Corollary~\ref{cor:uniform-characters}.
\end{theorem}

We work first with the candidate quotient
\[
   Q_\ell
   :=V^{-1}(D_\ell)/\langle\sigma(w_A)^r,\sigma(w_D)^2\rangle,
   \qquad
   A(Q_\ell)\cong U(D_\ell)/I_\ell^{\mathrm{cand}}.
\]
The identification \(Q_\ell\cong L_{-1}(D_\ell)\) from \cite{Jin2026}
is invoked only after its vacuum block has been classified.
Minimal Drinfeld--Sokolov reduction gives a lisse quotient whose Ramond Zhu
algebra \(B_\ell\) contains \(A_1\oplus D_r\) and satisfies
\begin{equation}
\label{eq:intro-nilpotence}
   e_A^r=0,\qquad e_D^2=0.
\end{equation}
Together with a weighted Casimir trace identity, these relations provide the
finite-dimensional bounds used below.

Set
\[
   J_{2,\ell}:=
   \operatorname{Ann}_{U(D_\ell)}L_{D_\ell}(-\alpha_2).
\]
All non-vacuum path tops have annihilator \(J_{2,\ell}\), so their descent
reduces to the inclusion
\begin{equation}
\label{eq:intro-membership}
   I_\ell^{\mathrm{cand}}\subset J_{2,\ell}.
\end{equation}
A double Li spectral flow and the Premet--Whittaker comparison prove
\eqref{eq:intro-membership}.  Independently, filtered Zhu theory reduces the
exhaustion problem to the zero and minimal nilpotent orbits; the finite-type
bounds and the type-\(D\) lattice gap then give
\begin{equation}
\label{eq:intro-exhaustion}
   \mu\in\{\mu_0,\mu_1,\ldots,\mu_\ell\}.
\end{equation}
These two statements classify the candidate quotient.  The
Grothendieck-group and character formulas then follow from the affine
Hecke-theoretic description of the subregular left cell.

Appendix~\ref{app:type-D-calculations} contains the type-\(D\) root
computations used in the finite-\(W\) comparison.  The rank-five case, where
\(D_3\cong A_3\), is recorded separately in
Appendix~\ref{app:rank-five-boundary}.

\section{Subregular affine cells and primitive ideals}
\label{sec:subregular-primitive}

Retain the notation of the introduction.  Let \(\widehat W\) be the affine
Weyl group and write
\(w\circ\lambda=w(\lambda+\widehat\rho)-\widehat\rho\).

\subsection{The distinguished subregular left cell}

For the shifted vacuum
\[
   A_\ell:=-\Lambda_0+\widehat\rho=N\Lambda_0+\rho
\]
one has
\[
   \langle A_\ell,\alpha_0^\vee\rangle=0,
   \qquad
   \langle A_\ell,\alpha_i^\vee\rangle=1\quad(1\le i\le\ell),
\]
so \(\operatorname{Stab}_{\widehat W}(A_\ell)=\langle s_0\rangle\).  The
distinguished Shan--Yan--Zhao element is therefore \(s_0\).  In the subregular
cell, unique reduced words ending in a fixed simple reflection form one left
cell \cite[Proposition~5.1]{BezrukavnikovKacKrylov2024}
\cite[Proposition~3.6]{Xu2019}.  Thus \(\mathbf c^L(s_0)\) consists of the
unique reduced paths in the affine Dynkin tree ending at node \(0\).

For \(1\le i\le\ell\), let \(z_i\) be the finite path word from node \(i\)
to node \(2\), ending in \(s_2\); explicitly
\begin{align*}
 z_1&=s_1s_2, & z_2&=s_2,\\
 z_i&=s_i s_{i-1}\cdots s_3s_2 &&(3\le i\le\ell-2),\\
 z_{\ell-1}&=s_{\ell-1}s_{\ell-2}\cdots s_3s_2,
 &z_\ell&=s_\ell s_{\ell-2}\cdots s_3s_2.
\end{align*}
Set \(w_0=s_0\) and \(w_i=z_i s_0\).

\begin{proposition}
\label{prop:subregular-left-cell}
\[
   \mathbf c^L(s_0)=\{w_0,w_1,\ldots,w_\ell\}.
\]
\end{proposition}

\begin{proof}
The affine \(D_\ell^{(1)}\)-diagram is a tree, hence there is exactly one
reduced path from each vertex to \(0\); these are precisely the words above.
\end{proof}

\begin{proposition}
\label{prop:path-dot-action}
Put \(\mu_0=0\) and \(\mu_i=z_i\rho-\rho\) for \(i\ge1\).  Then
\[
   w_i\circ(-\Lambda_0)=-\Lambda_0+\mu_i
   \qquad(0\le i\le\ell).
\]
\end{proposition}

\begin{proof}
Since \(s_0A_\ell=A_\ell\) and every \(z_i\) fixes \(\Lambda_0\),
\(z_i s_0A_\ell-\widehat\rho=-\Lambda_0+z_i\rho-\rho\).
\end{proof}

The path weights in orthogonal coordinates are
\begin{align}
 \mu_1&=(-2,1,1,0,\ldots,0),
 &\mu_2&=(0,-1,1,0,\ldots,0),\\
 \mu_i&=-\varepsilon_2-\cdots-\varepsilon_i+(i-1)\varepsilon_{i+1}
 &&(3\le i\le\ell-2),\\
 \mu_{\ell-1}&=-\varepsilon_2-\cdots-\varepsilon_{\ell-1}
                 +(\ell-2)\varepsilon_\ell,\\
 \mu_\ell&=-\varepsilon_2-\cdots-\varepsilon_{\ell-1}
                 -(\ell-2)\varepsilon_\ell.
\end{align}
These follow directly by applying the path reflections to \(\rho\).

\subsection{A common primitive annihilator}
\label{subsec:primitive-compression}

For a finite highest weight \(\lambda\), write
\[
   J(\lambda):=\operatorname{Ann}_{U(\mathfrak g)}L_{\mathfrak g}(\lambda),
   \qquad
   J_{2,\ell}:=J(-\alpha_2).
\]

\begin{theorem}
\label{thm:primitive-annihilator-compression}
For every \(1\le i\le\ell\),
\[
   J(\mu_i)=J_{2,\ell}.
\]
\end{theorem}

\begin{proof}
Each \(z_i\) is the unique reduced path ending in \(s_2\), hence
\(z_i\sim_Ls_2\) by \cite[Proposition~3.6]{Xu2019}.  The Duflo--Joseph
theorem for the regular integral block gives
\[
 \operatorname{Ann}L_{\mathfrak g}(x\rho-\rho)
 =\operatorname{Ann}L_{\mathfrak g}(y\rho-\rho)
 \iff x\sim_Ly
\]
(see \cite[Theorem~2.3.3(i)]{ShanYanZhao2026} for these conventions).
Since \(\mu_i=z_i\rho-\rho\) and \(s_2\rho-\rho=-\alpha_2\), this gives
\(J(\mu_i)=J_{2,\ell}\).
\end{proof}

\begin{corollary}
\label{cor:single-membership-test}
Let \(Q\) be a positive-energy quotient of \(V^{-1}(D_\ell)\) with
\(A(Q)\cong U(D_\ell)/I\).  For every \(1\le i\le\ell\),
\[
   L(-\Lambda_0+\mu_i)\text{ factors through }Q
   \quad\Longleftrightarrow\quad
   I\subset J_{2,\ell}.
\]
\end{corollary}

\begin{proof}
By Theorem~\ref{thm:primitive-annihilator-compression}, the assertion
reduces to Zhu descent.  If \(I\subset J_{2,\ell}\), the simple top
\(L_{\mathfrak g}(\mu_i)\) is an \(A(Q)\)-module.  Zhu correspondence
\cite[Theorem~4.9]{DongLiMason1998} gives its unique irreducible admissible
\(Q\)-lift; after inflation to \(V^{-1}(D_\ell)\), uniqueness for the
universal affine algebra identifies this lift with
\(L(-\Lambda_0+\mu_i)\).  The converse follows by taking the top.
\end{proof}

\begin{remark}
Direct minimal reduction cannot prove the non-vacuum descent.  Indeed,
\(z_i^{-1}\theta=s_2\theta=\theta-\alpha_2\), so
\(\langle\mu_i,\theta^\vee\rangle=-1\) and hence
\[
   \langle-\Lambda_0+\mu_i,\alpha_0^\vee\rangle=0.
\]
Arakawa's irreducible minimal-reduction theorem
\cite[Theorem~6.7.4]{Arakawa2005} therefore gives
\[
   H^0_{\mathrm{DS},f_\theta}\bigl(L(-\Lambda_0+\mu_i)\bigr)=0
   \qquad(1\le i\le\ell),
\]
which motivates the finite-\(W\) argument of Section~\ref{sec:membership}.
\end{remark}
\section{The candidate Zhu quotient and the minimal W-algebra}
\label{sec:candidate-zhu-minimalW}

Let \(\mathfrak g=D_\ell\), \(r=\ell-2\), \(N=2r+1\), and
\(K=r-1\).  In this section all constructions are made for the candidate
quotient; no simplicity statement for \(Q_\ell\) is used.

\subsection{The candidate quotient and its reduction}

Let \(w_A,w_D\) be highest vectors in the two Arakawa--Moreau summands of
\(S^2(\mathfrak g)\), and put
\[
   s_A=\sigma(w_A)^r,
   \qquad
   s_D=\sigma(w_D)^2.
\]
These are affine singular at level \(-1\)
\cite[Theorem~4.2]{ArakawaMoreau2018}.  Define
\begin{equation}
   Q_\ell:=V^{-1}(\mathfrak g)/\langle s_A,s_D\rangle,
   \qquad
   A(Q_\ell)\cong U(\mathfrak g)/I_\ell^{\mathrm{cand}},
\end{equation}
where \(I_\ell^{\mathrm{cand}}=\langle[s_A],[s_D]\rangle\).

For the minimal nilpotent \(f_\theta=e_{-\theta}\),
\[
   \mathfrak g^\natural\cong A_1\oplus D_r,
   \qquad
   \mathfrak g_{-1/2}
   \cong L_{A_1}(\varpi)\boxtimes L_{D_r}(\eta_1),
\]
and the two current levels are \(K=r-1\) and \(1\).  Let \(J_A,J_D\)
be highest-root currents.  Jin's \cite[Lemma~3.7 and Equation~(37)]{Jin2026}
place the defining sequence in the exactness category and show that the reduced
singular submodules are cyclic.  With \(c_A,c_D\in\mathbb C^\times\),
\cite[Lemma~3.3 and Proposition~3.6]{Jin2026} give
\[
   [s_A]_{\mathrm{DS}}=c_A:J_A^r:,
   \qquad
   [s_D]_{\mathrm{DS}}=c_D:J_D^2:.
\]
The generated-quotient lemma \cite[Proposition~2.3]{Jin2026}, or directly
\cite[Proposition~3.8]{Jin2026}, gives
\begin{equation}
\label{eq:exact-DS-Q}
   H^0_{\mathrm{DS},f_\theta}(Q_\ell)
   \cong
   \widetilde{\mathcal W}_\ell
   :=\mathcal W^{-1}(\mathfrak g,f_\theta)/
     \langle:J_A^r:,:J_D^2:\rangle.
\end{equation}
For the present level, \(Q_\ell\) is the quotient \(\widetilde V_{k_1}\)
with \(k_1=-1\) in \cite[Section~5]{ArakawaMoreau2018}.  Hence
\cite[Proposition~5.2]{ArakawaMoreau2018} gives
\(X_{Q_\ell}=\overline{\mathcal O_{\min}}\).  The associated-variety
formula for minimal reduction therefore shows that
\(\widetilde{\mathcal W}_\ell\) is nonzero and lisse.

Set
\[
   B_\ell:=A_{\mathrm R}(\widetilde{\mathcal W}_\ell).
\]
It is a nonzero finite-dimensional algebra
\cite[Section~3.3.1]{Jin2026}.  Writing
\(e_A=[J_A]_{\mathrm R}\) and \(e_D=[J_D]_{\mathrm R}\), the Ramond Zhu
product of mutually isotropic weight-one currents gives
\begin{equation}
   e_A^r=0,
   \qquad
   e_D^2=0.
\end{equation}

\begin{lemma}[Finite type bounds]
\label{lem:allowed-finite-types}
Every irreducible \(A_1\oplus D_r\)-constituent of a finite-dimensional
\(B_\ell\)-module is
\[
   L_{A_1}(a\varpi)\boxtimes L_{D_r}(\lambda_D),
\]
where
\begin{equation}
   0\le a\le r-1,
   \qquad
   \langle\lambda_D,\theta_D^\vee\rangle\le1.
\end{equation}
\end{lemma}

\begin{proof}
Let \(v\) be a highest vector of an irreducible constituent.  In the
highest-root \(\mathfrak{sl}_2\)-strings,
\(e_A^r f_A^r v\ne0\) if \(a\ge r\), while
\(e_D^2 f_D^2 v\ne0\) if
\(\langle\lambda_D,\theta_D^\vee\rangle\ge2\).  These contradict
\(e_A^r=0\) and \(e_D^2=0\), respectively.
\end{proof}

We normalize quadratic Casimirs by
\[
   c_2(a\varpi)=\frac{a(a+2)}2,
   \qquad
   c_2(\eta_1)=2r-1.
\]

\subsection{The weighted Casimir identity}

In the present Casimir normalization, Jin's level \(-1\) Ramond
relation \cite[Equation~(27)]{Jin2026} and contraction calculation
\cite[Lemma~3.11]{Jin2026} apply verbatim; in particular the mixed
\(A_1\)--\(D_r\) contribution has trace zero and the two diagonal
contraction coefficients are \(1\) and \(2/r\).

\begin{proposition}[Weighted Casimir trace identity]
Let \(M\ne0\) be a finite-dimensional \(B_\ell\)-module on which
\(L=[\omega]_{\mathrm R}\) acts by \(h\).  Then
\begin{equation}
\label{eq:weighted-Casimir}
   \frac{\operatorname{tr}_M\Omega_A}{\dim M}
   +
   \frac{r+2}{2r}
   \frac{\operatorname{tr}_M\Omega_D}{\dim M}
   =
   Nh+\frac{r-1}{4}.
\end{equation}
\end{proposition}

\begin{proof}
Put \(d=\dim M\) and \(\tau_*=\operatorname{tr}_M\Omega_*\).  Tracing
\cite[Equation~(27)]{Jin2026}, using
\cite[Lemma~3.11]{Jin2026}, and choosing
\(\langle u,v\rangle\ne0\) gives
\[
   2\tau_A+\left(1+\frac2r\right)\tau_D
   =2Nhd+\frac{r-1}{2}d.
\]
Dividing by \(2d\) yields \eqref{eq:weighted-Casimir}, in agreement with
\cite[Proposition~3.12]{Jin2026}.
\end{proof}

\subsection{Finite BRST and filtered Zhu}

Let \(H=U(\mathfrak g,f_\theta)\).  Applying Arakawa's Zhu--DS
compatibility for quotients \cite[Theorems~8.1 and~8.5]{Arakawa2015} to
\eqref{eq:exact-DS-Q} gives the finite realization needed later.

\begin{proposition}[Finite BRST realization of \(B_\ell\)]
\label{prop:B-finite-BRST}
\begin{equation}
   B_\ell
   \cong
   H^0_{f_\theta}
   \!\left(U(\mathfrak g)/I_\ell^{\mathrm{cand}}\right).
\end{equation}
\end{proposition}

\begin{proof}
Apply the twisted Zhu functor to \eqref{eq:exact-DS-Q} and use
\(A(Q_\ell)\cong U(\mathfrak g)/I_\ell^{\mathrm{cand}}\).
\end{proof}

We also need the following filtered-Zhu consequence for the exhaustion argument.
Equip \(U(\mathfrak g)\) with the PBW filtration and \(A(Q_\ell)\) with
Zhu's filtration.  Under the canonical affine identification
\[
   A(Q_\ell)\cong U(\mathfrak g)/I_\ell^{\mathrm{cand}},
\]
these are the same quotient filtration: the degree-one class
\(x(-1)\mathbf1\) maps to \(x\in\mathfrak g\).  Thus the standard Poisson
surjection
\(R_{Q_\ell}\twoheadrightarrow\operatorname{gr}A(Q_\ell)\)
\cite[Corollary~5.3]{ArakawaMoreau2018} is a surjection onto
\(\operatorname{gr}_{\mathrm{PBW}}(U(\mathfrak g)/I_\ell^{\mathrm{cand}})\).

\begin{lemma}[Filtered-Zhu containment]
\label{lem:filtered-Zhu-containment}
\begin{equation}
   \operatorname{Var}_{\mathrm{PBW}}
   \!\left(U(\mathfrak g)/I_\ell^{\mathrm{cand}}\right)
   \subset\overline{\mathcal O_{\min}}.
\end{equation}
Consequently, if \(I_\ell^{\mathrm{cand}}\subset J\) and \(J\) is
primitive, then
\[
   \mathcal V(J)=\{0\}
   \quad\text{or}\quad
   \mathcal V(J)=\overline{\mathcal O_{\min}}.
\]
\end{lemma}

\begin{proof}
Taking spectra of the preceding Poisson surjection gives the first
containment because
\(X_{Q_\ell}=\overline{\mathcal O_{\min}}\).  If
\(I_\ell^{\mathrm{cand}}\subset J\), the quotient map gives
\(\mathcal V(J)\subset
\operatorname{Var}_{\mathrm{PBW}}(U(\mathfrak g)/I_\ell^{\mathrm{cand}})\).
By Joseph's irreducibility theorem, the associated variety of a primitive
ideal is the closure of a single nilpotent orbit
\cite[Chapter~9]{Joseph1995}.  Since \(\mathcal O_{\min}\) is the unique
minimal nonzero nilpotent orbit, only the two displayed possibilities
remain.
\end{proof}
\section{Finite W-algebras and the membership theorem}
\label{sec:membership}

Retain the preceding notation, and let
\[
   H:=U(\mathfrak g,f_\theta)
\]
be the finite \(W\)-algebra associated with the minimal nilpotent.
By Proposition~\ref{prop:B-finite-BRST},
\begin{equation}
\label{eq:membership-B-finite-BRST}
   B_\ell
   \cong
   H^0_{f_\theta}
   \!\left(
      U(\mathfrak g)/I_\ell^{\mathrm{cand}}
   \right).
\end{equation}
We construct a simple \(B_\ell\)-module whose Skryabin lift has
annihilator \(J_{2,\ell}\).  By
Corollary~\ref{cor:single-membership-test}, this is enough to prove descent.

\subsection{A double Ramond spectral flow}

Recall that
\[
   \mathfrak g^\natural
   =
   \mathfrak s\oplus\mathfrak d,
   \qquad
   \mathfrak s\cong A_1,
   \qquad
   \mathfrak d\cong D_r,
\]
with current levels
\[
   k_A^\natural=K=r-1,
   \qquad
   k_D^\natural=1.
\]
We identify weights and coweights using the simply-laced normalization and
write
\[
   \varpi^\vee=\varpi=\frac{\theta_A}{2},
   \qquad
   \eta_1^\vee=\eta_1.
\]
Instead of the standard \(A_1\)-Ramond coweight \(\varpi^\vee\), consider
\begin{equation}
   x_*:=\varpi^\vee+\eta_1^\vee.
\end{equation}
Let
\[
   H_*:=J^{\{x_*\}}
   \in
   \widetilde{\mathcal W}_\ell{}_1.
\]
We apply Li's delta-operator construction to the adjoint module of
\(\widetilde{\mathcal W}_\ell\); see \cite{Li1996}.  Denote the resulting
twisted module by
\[
   \widetilde{\mathcal W}_\ell^{\,R,*}.
\]

\begin{lemma}
The inner automorphism
\[
   \exp(2\pi i\,x_{*,0})
\]
is the Ramond automorphism of the minimal \(W\)-algebra.  Equivalently, it
agrees with the automorphism obtained from the single coweight
\(\varpi^\vee\).
\end{lemma}

\begin{proof}
For every root \(\gamma\) of \(\mathfrak g^\natural\),
\[
   \gamma(\varpi^\vee)\in\mathbb Z,
   \qquad
   \gamma(\eta_1^\vee)\in\mathbb Z.
\]
Hence both current factors are fixed by the inner automorphism.

The weight-\(\frac32\) generators are parametrized by
\[
   U=\mathfrak g_{-\frac12}
   \cong
   L_{A_1}(\varpi)\boxtimes L_{D_r}(\eta_1).
\]
The two \(A_1\)-weights are \(\pm\varpi\), and
\[
   (\pm\varpi)(\varpi^\vee)\equiv\frac12\pmod{\mathbb Z}.
\]
Every weight of the \(D_r\)-vector representation belongs to
\(\{\pm\varepsilon_j\}\), and therefore has integral pairing with
\(\eta_1^\vee\).  Consequently every weight of \(U\) has
\[
   \nu(x_*)\equiv\frac12\pmod{\mathbb Z}.
\]
Thus the inner automorphism acts by \(-1\) on every
\(G^{\{u\}}\) and fixes the current fields and the conformal vector.  This is
precisely the Ramond automorphism.
\end{proof}

For a weight-one current \(J^{\{x\}}\), Li twisting gives
\[
   L_0^{R}
   =
   L_0+x_0+\frac{\kappa_x}{2},
\]
where \(\kappa_x\mathbf1=J^{\{x\}}_{(1)}J^{\{x\}}\).  Since the two current
factors are orthogonal,
\[
\begin{aligned}
   \kappa_{x_*}
   &=
   K(\varpi^\vee\mid\varpi^\vee)
   +(\eta_1^\vee\mid\eta_1^\vee)\\
   &=
   \frac{r-1}{2}+1
   =
   \frac{r+1}{2}.
\end{aligned}
\]
Hence
\begin{equation}
\label{eq:double-vacuum-energy}
   h_*=\frac{\kappa_{x_*}}2=\frac{r+1}{4}.
\end{equation}

\begin{lemma}[Lower boundedness]
The twisted adjoint module
\(\widetilde{\mathcal W}_\ell^{\,R,*}\) is lower bounded, and
\[
   \operatorname{Spec}L_0^{R,*}
   \subset
   h_*+\mathbb Z_{\ge0}.
\]
In particular, the vacuum vector \(\mathbf1\) belongs to the lowest Ramond
eigenspace.
\end{lemma}

\begin{proof}
The minimal \(W\)-algebra is strongly generated by the weight-one currents,
the fields \(G^{\{u\}}\) of conformal weight \(3/2\), and the conformal
vector.  The \(x_*\)-charges of current root vectors belong to
\[
   \{-1,0,1\}.
\]
The \(x_*\)-charges of the weights of \(U\) belong to
\[
   \left\{
      -\frac32,-\frac12,\frac12,\frac32
   \right\}.
\]
Thus every current generator has twisted excess
\[
   1+q\ge0,
\]
and every \(G\)-generator has twisted excess
\[
   \frac32+q\ge0.
\]
To make the passage from strong generators to arbitrary composite fields
explicit, for a homogeneous \(x_{*,0}\)-eigenvector \(a\) put
\[
   \epsilon_*(a):=\operatorname{wt}(a)+q_*(a).
\]
Since \(x_{*,0}\) is a derivation of the vertex algebra, charges add under
products.  Thus for homogeneous \(a,b\) and \(n\ge0\),
\begin{equation}
\label{eq:twisted-excess-additivity}
\begin{aligned}
   \operatorname{wt}\!\left(a_{(-n-1)}b\right)
      &=\operatorname{wt}(a)+\operatorname{wt}(b)+n,\\
   q_*\!\left(a_{(-n-1)}b\right)
      &=q_*(a)+q_*(b),
\end{aligned}
\end{equation}
and hence
\[
   \epsilon_*\!\left(a_{(-n-1)}b\right)
   =\epsilon_*(a)+\epsilon_*(b)+n\ge0.
\]
The conformal vector has excess \(2\), and derivatives increase excess by
positive integers.  Strong generation therefore implies that every state,
including all normally ordered products and their derivatives, has
nonnegative integral excess above the vacuum energy \(h_*\).  In particular,
no contraction term in a composite field can create a lower twisted
\(L_0\)-eigenvalue.
\end{proof}

\begin{remark}
Some current modes and some \(G\)-modes have zero twisted excess.  We
therefore do \emph{not} claim that the complete lowest Ramond eigenspace is
an irreducible \(\mathfrak g^\natural\)-module, or that it is generated only
by the vacuum under the \(A_1\)-current algebra.  The proof below uses only
the vacuum line inside that lowest eigenspace.
\end{remark}

The shifted Cartan zero modes are
\[
   J^{\{h\},R,*}_0
   =
   J^{\{h\}}_0+
   K(\varpi^\vee\mid h_A)
   +(\eta_1^\vee\mid h_D),
\]
for \(h=h_A+h_D\in\mathfrak h^\natural\).  Hence the vacuum has shifted
\(\mathfrak g^\natural\)-weight
\begin{equation}
\label{eq:double-vacuum-weight}
   \lambda_*
   =
   K\varpi+\eta_1
   =
   (r-1)\varpi+\eta_1.
\end{equation}

\subsection{Premet parameters of the twisted vacuum}
\label{subsec:Premet-vacuum-parameters}

We compare \eqref{eq:double-vacuum-weight} with Premet's highest-weight
coordinates for the minimal finite \(W\)-algebra \cite{Premet2007}.  Conjugate \(f_\theta\) to the minimal root vector
associated with
\[
   \beta=\alpha_2.
\]
We choose the compatible positive system for which the reductive part of the
centralizer is
\[
   A_1\oplus D_r.
\]

Premet's highest-weight theory involves a shift
\(\bar\delta\in(\mathfrak h^\natural)^*\).  The elementary type-\(D\) root
calculation is recorded in Appendix~\ref{app:type-D-calculations}.

\begin{lemma}[Premet shift]
\label{lem:Premet-shift-D}
In the \(\beta=\alpha_2\) model,
\begin{equation}
   \bar\delta
   =
   (r-1)\varpi+\eta_1.
\end{equation}
Consequently,
\[
   \lambda_*=\bar\delta.
\]
\end{lemma}

Let \(C\in Z(H)\) be Premet's quadratic central generator, normalized so that
the finite \(W\)-module associated with a Harish--Chandra parameter
\(\mu\) has eigenvalue \((\mu,\mu+2\rho)\).  Premet's type-\(D_\ell\)
calculation gives
\[
   c_0=-\ell(\ell-2)=-r(r+2).
\]
We record explicitly the normalization linking this generator to the Ramond
conformal class.  Put
\[
   L'=2(k+h^\vee)L+\frac12p_{\mathfrak g}(k).
\]
In the Ramond Zhu presentation, \cite[Equation~(5.1)]{KacMosenederFrajriaPapi2025}
writes the quadratic odd-generator relation in terms of
\(\Omega^\natural-L'\).  Under the identification with Premet's finite
\(W\)-algebra in \cite[Remark~5.8]{KacMosenederFrajriaPapi2025}, comparison
with \cite[Theorem~6.1(iv)]{Premet2007} identifies
\(L'=C-c_0\); the same remark records the compatible scalar check
\(c_0=-\frac12p_{\mathfrak g}(-h^\vee)\).
For \(D_\ell\), \cite[Section~3.1.2]{Jin2026} gives
\(p_{D_\ell}(k)=(k+2)(k+r)\).  Hence at \(k=-1\), since
\(k+h^\vee=2r+1=N\),
\begin{equation}
\label{eq:Premet-C-L-relation}
   C-c_0
   =
   2N L+\frac{r-1}{2}.
\end{equation}

\begin{lemma}
\label{lem:vacuum-C-zero}
The quadratic central element \(C\) acts on the twisted vacuum by zero.
\end{lemma}

\begin{proof}
By \eqref{eq:double-vacuum-energy} and
\eqref{eq:Premet-C-L-relation},
\[
\begin{aligned}
   C\mathbf1
   &=
   \left(
      -r(r+2)
      +2(2r+1)\frac{r+1}{4}
      +\frac{r-1}{2}
   \right)\mathbf1\\
   &=0.
\end{aligned}
\]
\end{proof}

Let
\[
   T_*:=
   \widetilde{\mathcal W}_\ell^{\,R,*}[h_*]
\]
be the lowest Ramond eigenspace.  It is naturally a module over
\(B_\ell\), hence, through \eqref{eq:membership-B-finite-BRST}, a module
over \(H\).

\begin{proposition}[The vacuum finite-\(W\) highest vector]
\label{prop:vacuum-finite-W-highest}
The vacuum vector
\[
   \mathbf1\in T_*
\]
is a finite-\(W\) highest-weight vector with Premet highest pair
\begin{equation}
   (\bar\delta,0).
\end{equation}
\end{proposition}

\begin{proof}
The two highest parameters are \(\bar\delta\) and \(0\) by
\eqref{eq:double-vacuum-weight}, Lemma~\ref{lem:Premet-shift-D}, and
Lemma~\ref{lem:vacuum-C-zero}.  Premet's positive part is generated by
\(\Theta(\mathfrak n^+(0))\) and \(\Theta(\mathfrak n^+(1))\)
\cite[Section~7]{Premet2007}.  Under the Ramond-Zhu/finite-\(W\)
identification \cite[Subsection~2.3]{Jin2026}, these give the positive
current and weight-\(3/2\) zero modes.

If \(a\) is an \(x_{*,0}\)-eigenvector of charge \(q\), Li's delta
operator gives, in the untwisted adjoint-mode notation,
\begin{equation}
\label{eq:Li-zero-mode-shift}
   a^{R,*}_0=a_q.
\end{equation}
For \(\Theta(\mathfrak n^+(0))\), Appendix~\ref{app:type-D-calculations}
gives charges \(q\in\{0,1\}\).  Premet's \(\mathfrak n^+(1)\) is spanned by
\(u_i^*=[e_\beta,z_i^*]\), of root \(\beta+\gamma_i^*\); since
\(\beta(x_*)=0\), its charge is \(\gamma_i^*(x_*)\), and the chosen positive
half has charges \(q\in\{\frac12,\frac32\}\).  Thus the positive finite-\(W\)
generators act on \(\mathbf1\) through the untwisted modes
\[
   J_q\quad(q=0,1),
   \qquad
   G_q\quad\left(q=\frac12,\frac32\right),
\]
all of which annihilate the vacuum by the creation property.  In particular,
the Li twist does not exchange Premet's chosen positive half with its dual.
Hence \(\mathbf1\) is a finite-\(W\) highest-weight vector.
\end{proof}

Let
\begin{equation}
   C_*:=B_\ell\mathbf1\subset T_*.
\end{equation}
Since \(B_\ell\) is finite dimensional, \(C_*\) is finite dimensional.  Any
proper submodule of \(C_*\) misses its cyclic generator \(\mathbf1\), so a
simple quotient of \(C_*\) retains a nonzero image of the vacuum highest
line.  Premet's highest-weight theory therefore gives the following.

\begin{corollary}
\label{cor:Mstar-B-module}
The finite-\(W\) simple module
\begin{equation}
   M_*:=L_H(\bar\delta,0)
\end{equation}
is a \(B_\ell\)-module.
\end{corollary}

\begin{proof}
Take a simple quotient of the cyclic \(B_\ell\)-module \(C_*\).  By
Proposition~\ref{prop:vacuum-finite-W-highest}, it is a simple highest-weight
\(H\)-module with highest pair \((\bar\delta,0)\).  The uniqueness of the
simple quotient of the corresponding finite-\(W\) Verma module
\cite[Section~7]{Premet2007} identifies it with
\(L_H(\bar\delta,0)\).
\end{proof}


\begin{lemma}[Rank-one Whittaker parameter comparison]
\label{lem:premet-ms-parameter}
Put \(\beta=\alpha_2\), and let \(\psi_\beta\) be the Whittaker
character supported on the \(\beta\)-root space, normalized by
\(\psi_\beta(e_\beta)=1\).  For \(\nu\in\mathfrak h^*\), write
\[
   \lambda=\bar\nu\in(\mathfrak h^\natural)^*,
   \qquad
   m=\langle\nu,\beta^\vee\rangle,
   \qquad
   c_\beta(\nu)=\frac12m(m+2).
\]
Here \(\bar\nu\) is the restriction to
\(\mathfrak h^\natural=\beta^\perp\), identified with the orthogonal
projection by Premet's bilinear form.  If \(M_{\rm P}(\lambda,c)\) denotes
Premet's parabolically induced module and \(M_{\rm MS}(\nu,\psi_\beta)\)
the Mili\v{c}i\'c--Soergel standard Whittaker module, then, after the
standard inner conjugation exchanging \(e_\beta\) and \(f_\beta\),
\begin{equation}
\label{eq:Premet-MS-standard-comparison}
   M_{\rm P}\!\left(\bar\nu,c_\beta(\nu)\right)
   \cong
   M_{\rm MS}(\nu,\psi_\beta).
\end{equation}
Moreover Premet's Whittaker functor gives
\begin{equation}
\label{eq:Premet-MS-finiteW-parameters}
   \operatorname{Wh}
   M_{\rm P}\!\left(\bar\nu,c_\beta(\nu)\right)
   \cong
   Z_H\!\left(
      \bar\nu+\bar\delta,
      (\nu,\nu+2\rho)
   \right).
\end{equation}
Both sides depend only on the
\(\langle s_\beta\rangle\)-dot orbit of \(\nu\).
\end{lemma}

\begin{proof}
For the character \(\psi_\beta\), the Mili\v{c}i\'c--Soergel support is
\(\Delta_{\psi_\beta}=\{\beta\}\).  Hence their parabolic and Levi are
\begin{equation}
\label{eq:MS-rank-one-parabolic}
   \mathfrak p_{\psi_\beta}
   =\mathfrak g_{\psi_\beta}\oplus\mathfrak n_\beta,
   \qquad
   \mathfrak g_{\psi_\beta}
   =\mathfrak h^\natural\oplus\mathfrak s_\beta,
   \qquad
   \mathfrak n_\beta
   =\bigoplus_{\gamma\in\Phi^+\setminus\{\beta\}}\mathfrak g_\gamma,
\end{equation}
with \(\mathfrak s_\beta\cong\mathfrak{sl}_2\); see
\cite[Section~2]{MilicicSoergel1997}.  This is precisely the parabolic
\(\mathfrak p_\beta\) used by Premet in the construction preceding
\cite[Theorem~1.3]{Premet2007}.

We compare the inducing modules themselves, rather than only their final
central characters.  In Premet's notation the inducing module is
\[
   Y_{\rm P}(\lambda,c)
   =U(\mathfrak p_\beta)/I_\beta(\lambda,c),
\]
where the left ideal \(I_\beta(\lambda,c)\) is generated by
\begin{equation}
\label{eq:Premet-inducing-relations}
   f_\beta-1,
   \qquad C_\beta-c,
   \qquad h-\lambda(h)\ \ (h\in\mathfrak h^\natural),
   \qquad e_\gamma\ \ (\gamma\in\Phi^+\setminus\{\beta\}),
\end{equation}
with
\[
   C_\beta=e_\beta f_\beta+f_\beta e_\beta
            +\frac12h_\beta^2;
\]
see the definition of \(M(\lambda,c)\) in
\cite[Section~7.3]{Premet2007}.  Mili\v{c}i\'c--Soergel define their
standard module by induction from the unique simple Whittaker module for
the Levi \(\mathfrak g_{\psi_\beta}\) with Harish--Chandra character
\(\xi_{\psi_\beta}(\nu)\), extending it trivially across
\(\mathfrak n_\beta\).

On the central \(\mathfrak h^\natural\)-factor this character is
\(\bar\nu\).  On the \(\mathfrak{sl}_2\)-factor, if
\(m=\langle\nu,\beta^\vee\rangle\), the Harish--Chandra value of
\(C_\beta\) is
\[
   m+\frac12m^2=\frac12m(m+2)=c_\beta(\nu).
\]
Thus the two Levi Whittaker modules have the same central data after
interchanging the two \(\beta\)-root spaces.  In the rank-one
\(\mathfrak{sl}_2\)-factor the nondegenerate simple Whittaker module with
prescribed central character is unique; together with the scalar
\(\mathfrak h^\natural\)-character, this identifies the two Levi inducing
modules once the two \(\beta\)-root spaces are exchanged.  Choose a representative
\(\dot s_\beta\) of \(s_\beta\), composed if necessary with a torus element,
so that
\[
   \operatorname{Ad}(\dot s_\beta)(f_\beta)=e_\beta.
\]
Because \(s_\beta\) fixes \(\mathfrak h^\natural=\beta^\perp\), preserves
\(C_\beta\), and permutes \(\Phi^+\setminus\{\beta\}\), the conjugate of
\eqref{eq:Premet-inducing-relations} is exactly the rank-one
Mili\v{c}i\'c--Soergel Whittaker inducing data with character
\(\psi_\beta(e_\beta)=1\) and Harish--Chandra parameter \(\nu\).  Therefore
\[
   {}^{\operatorname{Ad}(\dot s_\beta)}
   Y_{\rm P}\!\left(\bar\nu,c_\beta(\nu)\right)
   \cong
   Y_{\psi_\beta}\!\left(\xi_{\psi_\beta}(\nu),\psi_\beta\right).
\]
Inducing from the common parabolic \eqref{eq:MS-rank-one-parabolic} to
\(\mathfrak g\) proves \eqref{eq:Premet-MS-standard-comparison}.

The second finite-\(W\) parameter comes from Premet's Theorem~1.3.  For
\(\lambda\in(\mathfrak h^\natural)^*\),
\begin{equation}
\label{eq:Premet-Theorem13-explicit}
   \operatorname{Wh}M_{\rm P}(\lambda,c)
   \cong
   Z_H\!\left(
      \lambda+\bar\delta,
      c+(\lambda,\lambda+2\bar\rho)
   \right),
\end{equation}
where \(\bar\rho\) is the orthogonal projection of \(\rho\) to
\((\mathfrak h^\natural)^*\); see \cite[Theorem~1.3]{Premet2007}.
Since type \(D\) is simply laced and
\(\langle\rho,\beta^\vee\rangle=1\), the orthogonal decompositions are
\[
   \nu=\bar\nu+\frac m2\beta,
   \qquad
   \rho=\bar\rho+\frac12\beta.
\]
Consequently
\begin{align}
   (\nu,\nu+2\rho)
   &=
   (\bar\nu,\bar\nu+2\bar\rho)
   +\frac12m(m+2) \\
   &=
   (\bar\nu,\bar\nu+2\bar\rho)+c_\beta(\nu).
\end{align}
Substitution into \eqref{eq:Premet-Theorem13-explicit} proves
\eqref{eq:Premet-MS-finiteW-parameters}.
Finally \(s_\beta\mathbin\cdot\nu\) has the same restriction
\(\bar\nu\) and replaces \(m\) by \(-m-2\), leaving
\(c_\beta(\nu)\) unchanged.  This is also exactly the
Mili\v{c}i\'c--Soergel isomorphism criterion
\cite[Proposition~2.1]{MilicicSoergel1997}.
\end{proof}

\subsection{Identification of the primitive ideal $J_{2,\ell}$}

Recall from Section~\ref{subsec:primitive-compression} that
\[
   J_{2,\ell}
   =
   \operatorname{Ann}_{U(\mathfrak g)}
   L_{\mathfrak g}(-\alpha_2).
\]
We identify the Skryabin image of $M_*$ without choosing a
left/right-coset convention.

\begin{proposition}[Identification of $J_{2,\ell}$]
\label{prop:node2-identification}
Under Skryabin equivalence,
\begin{equation}
   \operatorname{Ann}_{U(\mathfrak g)}
   \operatorname{Skr}(M_*)
   =
   J_{2,\ell}.
\end{equation}
\end{proposition}

\begin{proof}
Put \(\beta=\alpha_2\).  Premet's Theorem~1.3 gives
\[
   \operatorname{Wh}M_{\rm P}(0,0)
   \cong Z_H(\bar\delta,0),
\]
and Whittaker--Skryabin is an exact equivalence.  Since
\(M_*=L_H(\bar\delta,0)\) is the unique simple quotient of this
finite-\(W\) Verma module, \(\operatorname{Skr}(M_*)\) is the unique
simple quotient of \(M_{\rm P}(0,0)\).

Lemma~\ref{lem:premet-ms-parameter}, applied to \(\nu=0\), identifies the
conjugate of \(M_{\rm P}(0,0)\) with
\(M_{\rm MS}(0,\psi_\beta)\).  Hence the conjugate of
\(\operatorname{Skr}(M_*)\) is
\(\mathcal L(0,\psi_\beta)\).  Since
\[
   s_\beta\mathbin\cdot0=-\beta,
\]
Proposition~2.1 of \cite{MilicicSoergel1997} gives
\[
   \mathcal L(0,\psi_\beta)
   \cong
   \mathcal L(-\beta,\psi_\beta).
\]
The weight \(-\beta\) is integral and
\(\langle-\beta+\rho,\beta^\vee\rangle=-1\), hence is
\(\langle s_\beta\rangle\)-antidominant.  Chen's annihilator theorem
\cite[Theorem~B]{Chen2021Annihilator} therefore yields
\[
   \operatorname{Ann}_{U(\mathfrak g)}
   \mathcal L(-\beta,\psi_\beta)
   =
   \operatorname{Ann}_{U(\mathfrak g)}L_{\mathfrak g}(-\beta)
   =J_{2,\ell}.
\]
The conjugation above is inner, and every two-sided ideal of
\(U(\mathfrak g)\) is invariant under inner automorphisms.  Therefore
\(\operatorname{Ann}\operatorname{Skr}(M_*)=J_{2,\ell}\), as claimed.
\end{proof}

\subsection{The membership theorem}

\begin{theorem}[Membership theorem]
\label{thm:membership}
For every \(\ell\ge5\),
\begin{equation}
\label{eq:membership}
   I_\ell^{\mathrm{cand}}
   \subset
   J_{2,\ell}.
\end{equation}
\end{theorem}

\begin{proof}
By Corollary~\ref{cor:Mstar-B-module} and
\eqref{eq:membership-B-finite-BRST},
\[
   M_*
   \in
   H^0_{f_\theta}
   \!\left(
      U(\mathfrak g)/I_\ell^{\mathrm{cand}}
   \right)\text{-Mod}.
\]
Arakawa's quotient Skryabin equivalence
\cite[Theorem~3.6]{Arakawa2015}, valid for arbitrary two-sided ideals, has
quasi-inverse given by the Skryabin functor; its Zhu realization is
\cite[Theorem~8.6]{Arakawa2015}.  Since \(I_\ell^{\mathrm{cand}}\) is the Zhu
ideal of \(Q_\ell\), \(\operatorname{Skr}(M_*)\) is annihilated by
\(I_\ell^{\mathrm{cand}}\).  Therefore
\[
   I_\ell^{\mathrm{cand}}\,
   \operatorname{Skr}(M_*)
   =
   0,
\]
and hence
\[
   I_\ell^{\mathrm{cand}}
   \subset
   \operatorname{Ann}_{U(\mathfrak g)}
   \operatorname{Skr}(M_*).
\]
Proposition~\ref{prop:node2-identification} identifies the annihilator on the
right with \(J_{2,\ell}\), proving
\eqref{eq:membership}.
\end{proof}

\begin{corollary}[Simultaneous descent]
For every \(1\le i\le\ell\), the affine highest-weight module
\[
   L(-\Lambda_0+\mu_i)
\]
factors through \(Q_\ell\).
\end{corollary}

\begin{proof}
By Theorem~\ref{thm:primitive-annihilator-compression},
\[
   \operatorname{Ann}_{U(\mathfrak g)}
   L_{\mathfrak g}(\mu_i)
   =
   J_{2,\ell}.
\]
Theorem~\ref{thm:membership} gives
\[
   I_\ell^{\mathrm{cand}}
   \subset
   J_{2,\ell}.
\]
Now apply Corollary~\ref{cor:single-membership-test}.
\end{proof}

\section{The exhaustion theorem}
\label{sec:exhaustion}

Retain the notation of the preceding sections, in particular the candidate
Zhu ideal \(I_\ell^{\mathrm{cand}}\subset U(\mathfrak g)\).

The affine vacuum-block congruence leaves no additional highest weights over
the candidate quotient:
\begin{equation}
\begin{gathered}
   I_\ell^{\mathrm{cand}}
   \subset
   \operatorname{Ann}_{U(\mathfrak g)}L_{\mathfrak g}(\mu),
   \qquad
   \mu+\rho\in W\rho+NQ(D_\ell),\\
   \Longrightarrow\qquad
   \mu\in\{\mu_0,\mu_1,\ldots,\mu_\ell\}.
\end{gathered}
\end{equation}

Write
\[
   J(\mu):=
   \operatorname{Ann}_{U(\mathfrak g)}L_{\mathfrak g}(\mu).
\]
By Lemma~\ref{lem:filtered-Zhu-containment},
\[
   I_\ell^{\mathrm{cand}}\subset J(\mu)
   \quad\Longrightarrow\quad
   \mathcal V(J(\mu))
   =
   \{0\}
   \quad\text{or}\quad
   \overline{\mathcal O_{\min}}.
\]
The two cases are handled separately.

\subsection{The zero-orbit branch}

Assume throughout this subsection that
\begin{equation}
\label{eq:zero-orbit-assumptions}
   I_\ell^{\mathrm{cand}}\subset J(\mu),
   \qquad
   \mathcal V(J(\mu))=\{0\},
   \qquad
   \mu+\rho\in W\rho+NQ(D_\ell).
\end{equation}
Since \(J(\mu)\) has zero associated variety,
\(L_{\mathfrak g}(\mu)\) is finite dimensional.  Hence
\[
   \mu\in P_+.
\]
Write its dominant orthogonal coordinates as
\[
   \mu=(x_1,\ldots,x_\ell),
   \qquad
   x_1\ge x_2\ge\cdots\ge x_{\ell-1}\ge |x_\ell|.
\]
The affine congruence in \eqref{eq:zero-orbit-assumptions} implies that all
\(x_i\) are integers, since both \(W\rho\) and \(NQ(D_\ell)\) have integral
orthogonal coordinates.

Consider the affine irreducible module
\[
   L(-\Lambda_0+\mu).
\]
Its affine simple-coroot value is
\[
   \left\langle
      -\Lambda_0+\mu,\alpha_0^\vee
   \right\rangle
   =
   -1-\langle\mu,\theta^\vee\rangle<0.
\]
Arakawa's irreducible minimal-reduction theorem therefore gives a nonzero
irreducible \(W\)-module
\begin{equation}
   \mathcal M_\mu
   :=
   H^0_{\mathrm{DS},f_\theta}
   \!\left(
      L(-\Lambda_0+\mu)
   \right)
   \ne0;
\end{equation}
see \cite[Theorem~6.7.4]{Arakawa2005}.  Since
\(I_\ell^{\mathrm{cand}}\subset J(\mu)\), the simple top
\(L_{\mathfrak g}(\mu)\) is an \(A(Q_\ell)\)-module; Zhu's unique
irreducible lift \cite[Theorem~4.9]{DongLiMason1998} therefore identifies
with \(L(-\Lambda_0+\mu)\).  Hence the affine module factors through
\(Q_\ell\), and \(\mathcal M_\mu\) factors through
\(\widetilde{\mathcal W}_\ell\).

Kac--Wakimoto's highest-weight formula identifies the
\(\mathfrak g^\natural\)-weight of the lowest Neveu--Schwarz space of
\(\mathcal M_\mu\) with the restriction of \(\mu\) to
\(\mathfrak h^\natural\)
\cite[Theorem~6.3]{KacWakimoto2004}.  We begin with the consequences of the Ramond Zhu relations.

\begin{lemma}[Restricted finite type]
\label{lem:zero-orbit-restricted-type}
Under \eqref{eq:zero-orbit-assumptions},
\[
   \mu
   =
   (x_1,x_2,c,0,\ldots,0),
   \qquad
   c\in\{0,1\}.
\]
If
\[
   a:=x_1-x_2,
\]
then
\[
   0\le a\le r-1.
\]
Equivalently,
\[
   \mu|_{\mathfrak h^\natural}
   =
   a\varpi+c\eta_1.
\]
\end{lemma}

\begin{proof}
The lowest Neveu--Schwarz \(\mathfrak g^\natural\)-module is finite
dimensional because \(L_{\mathfrak g}(\mu)\) is finite dimensional.  Let
\(u\) be a vector in this lowest space which is lowest weight for the
\(A_1\)-factor and highest weight for the \(D_r\)-factor.

Twist \(\mathcal M_\mu\) by the single coweight
\[
   x=\varpi^\vee.
\]
The \(A_1\)-current charges are in \(\{-1,0,1\}\), while the
weight-\(\frac32\) generators have \(x\)-charges \(\pm\frac12\).
For a homogeneous \(x_0\)-eigenvector \(a\), write
\[
   \epsilon_x(a):=\operatorname{wt}(a)+q_x(a).
\]
Thus the current generators have excess \(1+q_x\ge0\), the
weight-\(\frac32\) generators have excess
\(\frac32+q_x\ge1\), and the conformal vector has excess \(2\).
Moreover, since \(x_0\) is a derivation, for homogeneous \(a,b\) and
\(n\ge0\),
\[
   \epsilon_x\!\left(a_{(-n-1)}b\right)
   =\epsilon_x(a)+\epsilon_x(b)+n.
\]
Strong generation by the currents, the \(G\)-fields, and the conformal
vector therefore implies that every creation monomial has nonnegative
twisted excess.  Since \(\mathcal M_\mu\) is generated from its lowest
Neveu--Schwarz space by such creation modes, and \(u\) has the smallest
\(x\)-weight in that ground space, no vector can have twisted conformal
weight strictly below that of \(u\).  Hence \(u\) lies in the lowest
Ramond eigenspace.  No irreducibility statement about the complete
Ramond bottom is being used.

The cyclic module
\[
   C_\mu:=B_\ell u
\]
is finite dimensional.  Under the twisted \(A_1\)-zero modes,
\(u\) is a highest-weight vector of highest weight
\[
   (K-a)\varpi.
\]
Finite-dimensionality therefore forces
\[
   K-a\ge0,
\]
or \(a\le r-1\).

The \(D_r\)-factor is unchanged by this twist, and \(u\) remains a
\(D_r\)-highest vector.  The relation \(e_D^2=0\), equivalently
Lemma~\ref{lem:allowed-finite-types}, gives
\[
   \left\langle
      \lambda_D,\theta_D^\vee
   \right\rangle
   \le1
\]
for its \(D_r\)-highest weight \(\lambda_D\).  The level-one dominant
weights are
\[
   0,\qquad
   \eta_1,\qquad
   \eta_{r-1},\qquad
   \eta_r
\]
for \(r\ge4\), with the analogous four weights for
\(D_3\cong A_3\).  The affine congruence makes the orthogonal coordinates of
\(\mu\) integral, whereas the two spinor weights have half-integral
orthogonal coordinates.  Thus
\[
   \lambda_D=c\eta_1,
   \qquad
   c\in\{0,1\}.
\]
Dominance of \(\mu\) now gives
\[
   \mu=(x_1,x_2,c,0,\ldots,0).
\]
\end{proof}

Set
\begin{equation}
   a:=x_1-x_2,
   \qquad
   s:=x_1+x_2
   =
   \langle\mu,\theta^\vee\rangle.
\end{equation}

\begin{lemma}[Lowest conformal energies]
\label{lem:zero-orbit-energies}
The lowest Neveu--Schwarz conformal weight of \(\mathcal M_\mu\) is
\begin{equation}
\label{eq:zero-hW}
   h_W
   =
   \frac{
      s^2+a(a+2)+2c(2r-1)
   }{4N}.
\end{equation}
For the single Li twist \(x=\varpi^\vee\), the Ramond ground vector \(u\)
above has energy
\begin{equation}
\label{eq:zero-hR}
   h_R
   =
   h_W-\frac a2+\frac{r-1}{4}.
\end{equation}
\end{lemma}

\begin{proof}
Kac--Wakimoto's conformal-weight formula
\cite[Equation~(7.6)]{KacWakimoto2004} reads, in the present normalization,
\[
   h_W
   =
   \frac{
      \bigl((\mu+\rho,\theta)-N\bigr)^2-(k+1)^2
      +2(\mu^\natural,\mu^\natural+2\rho^\natural)
   }{4N}.
\]
At \(k=-1\),
\[
   (\rho,\theta)=N,
   \qquad
   k+1=0.
\]
Moreover,
\[
   \mu^\natural=a\varpi+c\eta_1,
\]
and, since \(c\in\{0,1\}\),
\[
   (a\varpi,a\varpi+2\rho_{A_1})
   =\frac{a(a+2)}2,
   \qquad
   (c\eta_1,c\eta_1+2\rho_{D_r})
   =c(2r-1).
\]
This gives \eqref{eq:zero-hW}.

The vector \(u\) has \(x\)-weight \(-a/2\).  For Li's operator
\cite[Proposition~5.4]{Li1996},
\[
   \Delta_x(z)\omega
   =\omega+J^{\{x\}}z^{-1}+\frac{\kappa_x}{2}z^{-2},
   \qquad
   \kappa_x=K(\varpi^\vee\mid\varpi^\vee)=\frac K2.
\]
Hence
\[
   L_0^R=L_0+x_0+\frac{\kappa_x}{2}
   =L_0+x_0+\frac K4,
\]
in agreement with \cite[Equation~(62)]{Jin2026}; this yields
\eqref{eq:zero-hR}.
\end{proof}

The cyclic Ramond ground module \(C_\mu=B_\ell u\) may contain several
\(\mathfrak g^\natural\)-constituents.  We therefore use a uniform Casimir
upper bound rather than attempting to identify the complete bottom space.

\begin{lemma}[Root--vector coset bound]
\label{lem:zero-root-vector-bound}
Every \(D_r\)-constituent of \(C_\mu\) belongs to the root or vector coset of
the \(D_r\)-weight lattice.  Consequently
\begin{equation}
\label{eq:zero-Cmax}
   \frac{\operatorname{tr}_{C_\mu}\Omega_A}{\dim C_\mu}
   +
   \frac{r+2}{2r}
   \frac{\operatorname{tr}_{C_\mu}\Omega_D}{\dim C_\mu}
   \le
   C_{\max},
\end{equation}
where
\begin{equation}
\label{eq:Cmax-definition}
   C_{\max}
   :=
   \frac{r^2-1}{2}
   +
   \frac{r+2}{2r}(2r-1).
\end{equation}
\end{lemma}

\begin{proof}
The initial \(D_r\)-highest weight is \(c\eta_1\), hence belongs to the
subgroup of \(P(D_r)/Q(D_r)\) generated by the vector class.  Current
zero modes change weights by roots, while the non-current generators
\(G^{\{u\}}\) carry \(D_r\)-weights in
\[
   \eta_1+Q(D_r).
\]
Since \(\widetilde{\mathcal W}_\ell\) is strongly generated by the currents,
the \(G\)-fields, and the conformal vector, its Ramond Zhu algebra
\(B_\ell\) is generated by their Zhu classes.  Thus every operator used to
generate \(C_\mu=B_\ell u\) changes the \(D_r\)-coset by either the root
class or the vector class (the conformal class has weight zero).
Therefore the cyclic module can meet only the root and vector cosets, never
a spinor coset.  This remains true for \(r=3\cong A_3\), where the subgroup
generated by the vector weight \(\omega_2\) is
\(\{0,\omega_2\}\) modulo the root lattice.

By Lemma~\ref{lem:allowed-finite-types}, the largest \(A_1\)-Casimir among
allowed constituents is attained at \((r-1)\varpi\) and equals
\[
   \frac{(r-1)(r+1)}2
   =
   \frac{r^2-1}{2}.
\]
Among the allowed root/vector \(D_r\)-types, the largest Casimir is the
vector value
\[
   c_2(\eta_1)=2r-1.
\]
Averaging over irreducible constituents gives \eqref{eq:zero-Cmax}.
\end{proof}

\begin{proposition}[Zero-orbit energy gap]
Under \eqref{eq:zero-orbit-assumptions},
\begin{equation}
   s<N.
\end{equation}
\end{proposition}

\begin{proof}
Apply the weighted Casimir identity
\eqref{eq:weighted-Casimir} to \(C_\mu\):
\[
   Nh_R+\frac{r-1}{4}
   =
   \frac{\operatorname{tr}_{C_\mu}\Omega_A}{\dim C_\mu}
   +
   \frac{r+2}{2r}
   \frac{\operatorname{tr}_{C_\mu}\Omega_D}{\dim C_\mu}.
\]
Combining Lemmas~\ref{lem:zero-orbit-energies} and
\ref{lem:zero-root-vector-bound} gives
\begin{equation}
   s^2
   \le
   4ar-a^2
   +4r+6-\frac4r
   -2c(2r-1).
\end{equation}
Since
\[
   0\le a\le r-1,
   \qquad
   c\ge0,
\]
the right-hand side is at most
\[
   3r^2+2r+5-\frac4r.
\]
But
\[
\begin{aligned}
   N^2-
   \left(
      3r^2+2r+5-\frac4r
   \right)
   &=
   (2r+1)^2-
   \left(
      3r^2+2r+5-\frac4r
   \right)\\
   &=
   r^2+2r-4+\frac4r\\
   &>0
\end{aligned}
\]
for every \(r\ge3\).  Since \(s\ge0\), this proves \(s<N\).
\end{proof}

Now impose the affine congruence.  Recall
\[
   Q(D_\ell)
   =
   \left\{
      (q_1,\ldots,q_\ell)\in\mathbb Z^\ell
      \,\middle|\,
      \sum_iq_i\equiv0\pmod2
   \right\},
\]
and
\[
   W(D_\ell)
\]
is the group of signed permutations with an even number of sign changes.

\begin{proposition}[Zero-orbit rigidity]
\label{prop:zero-orbit-rigidity}
Under \eqref{eq:zero-orbit-assumptions},
\[
   \mu=0.
\]
\end{proposition}

\begin{proof}
By Lemma~\ref{lem:zero-orbit-restricted-type},
\[
   \mu+\rho
   =
   (r+1+x_1,\ r+x_2,\ r-1+c,\ r-2,\ldots,1,0).
\]
Reduce the affine congruence modulo
\[
   N=2r+1.
\]
The final coordinates
\[
   r-2,r-3,\ldots,1,0
\]
are all strictly smaller than \(r\).  A negative signed coordinate
\(-t\), with \(1\le t\le r+1\), has residue
\[
   N-t\ge r.
\]
Hence these final coordinates can only arise from the corresponding
positive entries of \(W\rho\).  They consume the absolute values
\[
   0,1,\ldots,r-2.
\]
Only
\[
   r-1,\quad r,\quad r+1
\]
remain for the first three coordinates.

Suppose first that \(c=0\).  The third residue is \(r-1\), which can only
come from the positive entry \(r-1\).  Thus the first two residues must be
those obtained from the signed entries \(r\) and \(r+1\), and hence each is
\(r\) or \(r+1\).  Since \(s<N\),
\[
   x_2\le r,
\]
so
\[
   r+x_2\in\{r,r+1\}
\]
forces \(x_2\in\{0,1\}\).  Also
\[
   x_1-x_2\le r-1,
\]
hence \(x_1\le r\).  If \(x_1=r\), the first coordinate has residue zero,
which is impossible; therefore
\[
   r+1+x_1<N.
\]
Its residue must be \(r\) or \(r+1\), so necessarily \(x_1=0\).  Dominance
then gives \(x_2=0\).

Now suppose that \(c=1\).  The third residue is \(r\).  It can arise only
from either \(+r\) or \(-(r+1)\).  In both cases, examination of the two
remaining absolute values shows that
\[
   r+x_2\in\{r+1,r+2\},
\]
hence
\[
   x_2\in\{1,2\}.
\]
If \(x_2=2\), then the residue \(r+2\) must come from \(-(r-1)\),
so the first residue is \(r\) or \(r+1\).  But
\(2\le x_1\le r+1\), by \(x_1-x_2\le r-1\), and hence
\(r+1+x_1\pmod N\) lies in
\(\{r+3,\ldots,2r,0,1\}\), a contradiction.  Thus \(x_2=1\); the
remaining residue condition then gives \(r+1+x_1=r+2\), so \(x_1=1\).

For \((x_1,x_2,c)=(1,1,1)\), the first three entries of \(w\rho\)
can only be
\[
   (-(r-1),r+1,r)\quad\text{or}\quad (-(r-1),-r,-(r+1)).
\]
The sign parity does not exclude these, since the zero coordinate of
\(\rho\) can absorb an extra sign change.  Instead set
\(q=(\mu+\rho-w\rho)/N\).  The two cases give respectively
\[
   q=(1,0,0,\ldots,0),\qquad q=(1,1,1,0,\ldots,0),
\]
both of odd coordinate sum.  Hence \(q\notin Q(D_\ell)\), contradicting
the affine congruence.  Thus \(c=1\) is impossible.

We conclude
\[
   c=0,\qquad x_1=x_2=0,
\]
and therefore \(\mu=0\).
\end{proof}

\subsection{The minimal-orbit branch}

Assume now
\begin{equation}
\label{eq:minimal-orbit-assumptions}
   I_\ell^{\mathrm{cand}}\subset J(\mu),
   \qquad
   \mathcal V(J(\mu))
   =
   \overline{\mathcal O_{\min}},
   \qquad
   \xi:=\mu+\rho\in W\rho+NQ(D_\ell).
\end{equation}

Since \(\mathcal V(J(\mu))=\overline{\mathcal O_{\min}}\) contains
\(f_\theta\), the finite-\(W\)/primitive-ideal correspondence
\cite[Proposition~3.1]{Petukhov2018} gives a finite-dimensional simple module
\(M_J\) over
\[
   H=U(\mathfrak g,f_\theta)
\]
such that
\begin{equation}
   \operatorname{Ann}_{U(\mathfrak g)}
   \operatorname{Skr}(M_J)
   =
   J(\mu);
\end{equation}
Since \(\operatorname{Skr}(M_J)\) is annihilated by \(J(\mu)\), it is
annihilated by \(I_\ell^{\mathrm{cand}}\).  Arakawa's quotient Skryabin
 equivalence \cite[Theorem~3.6]{Arakawa2015}, together with
\eqref{eq:membership-B-finite-BRST}, therefore makes
\(M_J=\operatorname{Wh}(\operatorname{Skr}(M_J))\) a finite-dimensional
\(B_\ell\)-module.

Let \(C\in Z(H)\) be the quadratic central element normalized as in
Section~\ref{subsec:Premet-vacuum-parameters}.  It acts on \(M_J\) by
\begin{equation}
\label{eq:C-HC-eigenvalue}
   C
   =
   (\mu,\mu+2\rho)
   =
   \|\xi\|^2-\|\rho\|^2.
\end{equation}

The next lemma relates a primitive ideal to the highest-weight parameters
of the corresponding simple finite \(W\)-module by identifying the
Mili\v{c}i\'c--Soergel Whittaker model with Premet's highest-weight model.

\begin{lemma}[Primitive ideals and finite \(W\) highest weights]
\label{lem:primitive-finiteW-bridge}
Let \(J\subset U(\mathfrak g)\) be a primitive ideal with
\[
   \mathcal V(J)=\overline{\mathcal O_{\min}}
\]
and integral central character, and let \(M_J\) be the finite-dimensional
simple \(H\)-module characterized by
\[
   \operatorname{Ann}_{U(\mathfrak g)}\operatorname{Skr}(M_J)=J.
\]
Put \(\beta=\alpha_2\).  Then there exists an integral
\(\langle s_\beta\rangle\)-antidominant weight \(\nu\) such that
\[
   \operatorname{Ann}_{U(\mathfrak g)}L_{\mathfrak g}(\nu)=J,
\]
and, in Premet's \(\beta\)-model,
\begin{equation}
\label{eq:primitive-finiteW-highest-bridge}
   M_J
   \cong
   L_H\!\left(
      \overline{\nu}+\overline{\delta},
      (\nu,\nu+2\rho)
   \right).
\end{equation}
In particular, the \(\mathfrak g^\natural\)-highest weight of \(M_J\) is
\[
   \lambda_J^H=\overline{\nu}+\overline{\delta}.
\]
If a prescribed \(\nu_0\) with
\(\operatorname{Ann}L_{\mathfrak g}(\nu_0)=J\) is already
\(\langle s_\beta\rangle\)-antidominant, then one may take
\(\nu=\nu_0\).
\end{lemma}

\begin{proof}
Write the Premet highest pair of \(M_J\) as
\[
   (\lambda_J^H,C_J).
\]
Set
\[
   \lambda=\lambda_J^H-\bar\delta,
   \qquad
   c=C_J-(\lambda,\lambda+2\bar\rho).
\]
By Premet's Theorem~1.3,
\[
   \operatorname{Wh}M_{\rm P}(\lambda,c)
   \cong
   Z_H(\lambda_J^H,C_J).
\]
Since \(M_J\) is the unique simple quotient of this finite-\(W\) Verma
module and Whittaker--Skryabin is an exact equivalence,
\(\operatorname{Skr}(M_J)\) is the unique simple quotient of
\(M_{\rm P}(\lambda,c)\).

Choose \(m\in\mathbb C\) with
\[
   \frac12m(m+2)=c
\]
and let \(\nu\in\mathfrak h^*\) be the extension of \(\lambda\) determined
by
\[
   \bar\nu=\lambda,
   \qquad
   \langle\nu,\beta^\vee\rangle=m.
\]
The two choices of \(m\) are \(m\) and \(-m-2\), hence the two resulting
weights form one \(\langle s_\beta\rangle\)-dot orbit.  By
Lemma~\ref{lem:premet-ms-parameter}, after the fixed inner conjugation the
module \(M_{\rm P}(\lambda,c)\) is the Mili\v{c}i\'c--Soergel standard
module \(M_{\rm MS}(\nu,\psi_\beta)\).  Replace \(\nu\), if necessary,
by the \(\langle s_\beta\rangle\)-antidominant representative of this
orbit.  Proposition~2.1 of \cite{MilicicSoergel1997} shows that neither the
standard module nor its simple quotient changes.  Consequently the simple
quotient \(\mathcal L(\nu,\psi_\beta)\) has two-sided annihilator \(J\).

The full central character of \(M_{\rm MS}(\nu,\psi_\beta)\) is
\(\xi(\nu)\); see \cite[Section~2]{MilicicSoergel1997}.  Since the central
character of \(J\) is integral, there is an integral weight \(\lambda_0\)
with the same Harish--Chandra character.  Hence
\[
   \nu=w\mathbin\cdot\lambda_0
\]
for some \(w\in W\), and therefore \(\nu\) is integral because the dot
action preserves the weight lattice.  By construction \(\nu\) is also the
\(W_{\psi_\beta}=\langle s_\beta\rangle\)-antidominant representative, so
it satisfies the parameter condition in Chen's annihilator theorem.
That theorem now applies and gives
\[
   \operatorname{Ann}_{U(\mathfrak g)}L_{\mathfrak g}(\nu)
   =
   \operatorname{Ann}_{U(\mathfrak g)}
      \mathcal L(\nu,\psi_\beta)
   =J.
\]
Finally Lemma~\ref{lem:premet-ms-parameter} identifies the finite-\(W\)
highest pair of the simple quotient with
\[
   \left(
      \bar\nu+\bar\delta,
      (\nu,\nu+2\rho)
   \right).
\]
By uniqueness of the simple quotient of a finite-\(W\) Verma module this
is the highest pair of \(M_J\), proving
\eqref{eq:primitive-finiteW-highest-bridge}.

Now suppose that a prescribed \(\nu_0\) is already
\(\langle s_\beta\rangle\)-antidominant and
\(\operatorname{Ann}L_{\mathfrak g}(\nu_0)=J\).  Chen's theorem gives
\[
   \operatorname{Ann}\mathcal L(\nu_0,\psi_\beta)=J.
\]
Let \(V_0\) be its inverse image under Whittaker--Skryabin.  Premet's
bijection of primitive spectra \cite[Theorem~1.2(i)]{Premet2007} implies
that \(\operatorname{Ann}_H V_0=\operatorname{Ann}_H M_J\).  Since
\(M_J\) is finite dimensional, the quotient by this common annihilator is
a full matrix algebra; hence it has a unique simple module.  Thus
\(V_0\cong M_J\).  Applying Lemma~\ref{lem:premet-ms-parameter} to
\(\nu_0\) gives exactly the pair
\[
   \left(
      \bar\nu_0+\bar\delta,
      (\nu_0,\nu_0+2\rho)
   \right),
\]
so \(\nu_0\) may be used in
\eqref{eq:primitive-finiteW-highest-bridge}.
\end{proof}

\begin{lemma}[No spinor cosets]
\label{lem:minimal-no-spinor}
Every \(D_r\)-constituent of \(M_J\) belongs to the root or vector coset.
\end{lemma}

\begin{proof}
Apply Lemma~\ref{lem:primitive-finiteW-bridge} to \(J=J(\mu)\), and let
\(\nu\) be the resulting \(\langle s_2\rangle\)-antidominant weight.  Since
\(J(\mu)=\operatorname{Ann}L_{\mathfrak g}(\nu)\), the two highest-weight
modules have the same central character; hence
\[
   \nu+\rho=w(\mu+\rho)=w\xi
\]
for some \(w\in W\).  The set \(W\rho+NQ(D_\ell)\) is \(W\)-stable, so
\[
   \nu+\rho\in W\rho+NQ(D_\ell).
\]
In particular, \(\nu+\rho\) has integral orthogonal coordinates.  By
Lemma~\ref{lem:primitive-finiteW-bridge},
\begin{equation}
   \lambda_H
   =
   \overline{\nu}+\overline{\delta}
   =
   \overline{\nu+\rho}
   -
   \overline{\rho}
   +
   \overline{\delta}.
\end{equation}
In the \(D_r\)-factor, \(\overline{\rho}\) and
\(\overline{\delta}\) have integral orthogonal coordinates.  Hence the
\(D_r\)-component of \(\lambda_H\) has integral orthogonal coordinates and
cannot lie in either spinor coset.

As in Lemma~\ref{lem:zero-root-vector-bound}, the current generators move
weights by roots and the \(G\)-generators move them by vector-coset weights.
Thus the entire cyclic finite-\(W\) module, and hence every
\(D_r\)-constituent of the simple module \(M_J\), stays in the root/vector
subgroup of \(P(D_r)/Q(D_r)\).
\end{proof}

\begin{proposition}[Quadratic central gap]
\label{prop:minimal-Casimir-gap}
Under \eqref{eq:minimal-orbit-assumptions},
\begin{equation}
   C\le 2-\frac2r<2.
\end{equation}
\end{proposition}

\begin{proof}
Since \(M_J\) is finite dimensional over \(B_\ell\), the finite-type
bound of Lemma~\ref{lem:allowed-finite-types}, together with
Lemma~\ref{lem:minimal-no-spinor}, gives
\[
   \overline{\mathcal C}(M_J)
   :=
   \frac{\operatorname{tr}_{M_J}\Omega_A}{\dim M_J}
   +
   \frac{r+2}{2r}
   \frac{\operatorname{tr}_{M_J}\Omega_D}{\dim M_J}
   \le C_{\max},
\]
where \(C_{\max}\) is defined in \eqref{eq:Cmax-definition}.

The relation \eqref{eq:Premet-C-L-relation} and the weighted trace identity
\eqref{eq:weighted-Casimir} give
\[
   C
   =
   2\overline{\mathcal C}(M_J)-r(r+2).
\]
Therefore
\[
\begin{aligned}
   C
   &\le
   2C_{\max}-r(r+2)\\
   &=
   2-\frac2r,
\end{aligned}
\]
which is strictly smaller than \(2\) for \(r\ge3\).
\end{proof}

The affine lattice gives a much larger gap away from the regular central
character.

\begin{lemma}[Type-\(D\) norm gap]
\label{lem:D-norm-gap}
Let
\[
   \zeta\in W\rho+NQ(D_\ell).
\]
Then
\[
   \|\zeta\|^2-\|\rho\|^2\ge0.
\]
Moreover,
\[
   \|\zeta\|^2-\|\rho\|^2=0
   \quad\Longleftrightarrow\quad
   \zeta\in W\rho,
\]
while otherwise
\begin{equation}
\label{eq:D-norm-positive-gap}
   \|\zeta\|^2-\|\rho\|^2
   \ge2N.
\end{equation}
\end{lemma}

\begin{proof}
After applying an element of \(W\), write
\[
   \zeta=\rho+Nq,
   \qquad
   q\in Q(D_\ell).
\]
Then
\begin{equation}
\label{eq:D-norm-Rq}
   \|\zeta\|^2-\|\rho\|^2
   =
   N R(q),
   \qquad
   R(q):=
   N\|q\|^2+2(\rho,q).
\end{equation}
Write
\[
   \rho=(r+1,r,r-1,\ldots,1,0).
\]
For a coordinate of \(\rho\) equal to \(t\), put
\[
   f_t(n)=Nn^2+2tn,
   \qquad n\in\mathbb Z.
\]
Since \(N=2r+1\), one has \(f_t(n)\ge0\) for
\(0\le t\le r\), with equality only at \(n=0\).  For the first
coordinate, \(t=r+1\), the only negative value is
\[
   f_{r+1}(-1)=-1;
\]
all other nonzero values are positive.  Hence, if \(q_1\ne-1\), every
coordinate contribution is nonnegative and \(R(q)=0\) forces \(q=0\).

Suppose instead that \(q_1=-1\).  Since
\[
   Q(D_\ell)=\{(q_i)\in\mathbb Z^\ell:\ \sum_iq_i\in2\mathbb Z\},
\]
the remaining coordinates have odd total sum, so at least one of them is
odd.  Among all odd \(n\) and all \(0\le t\le r\), the smallest value of
\(f_t(n)\) is
\[
   f_r(-1)=1,
\]
and this minimum is unique.  Thus the contribution \(-1\) from the first
coordinate is compensated by at least \(1\), with equality only when
\[
   q=(-1,-1,0,\ldots,0)=-\theta.
\]
Consequently
\[
   R(q)\ge0,
\]
with equality precisely for
\[
   q=0
   \qquad\text{or}\qquad
   q=-\theta.
\]
In the latter case
\[
   \rho-N\theta=s_\theta\rho
\]
because
\[
   \langle\rho,\theta^\vee\rangle=N.
\]
Thus equality in \eqref{eq:D-norm-Rq} is equivalent to
\(\zeta\in W\rho\).

Finally, \(Q(D_\ell)\) is an even lattice, so
\[
   \|q\|^2\in2\mathbb Z.
\]
Hence \(R(q)\in2\mathbb Z\).  If \(R(q)>0\), then \(R(q)\ge2\), which gives
\eqref{eq:D-norm-positive-gap}.  The bound is sharp: for
\(q=(-1,0,-1,0,\ldots,0)\) one has \(R(q)=2\).
\end{proof}

\begin{proposition}[Regularity of the minimal-orbit central character]
\label{prop:minimal-regular-central-character}
Under \eqref{eq:minimal-orbit-assumptions},
\[
   \xi=\mu+\rho\in W\rho.
\]
Equivalently, \(J(\mu)\) has the regular integral central character
\(\mathfrak m_0\).
\end{proposition}

\begin{proof}
By \eqref{eq:C-HC-eigenvalue} and Lemma~\ref{lem:D-norm-gap},
either
\[
   C=0
\]
and \(\xi\in W\rho\), or
\[
   C\ge2N>2.
\]
The second possibility contradicts
Proposition~\ref{prop:minimal-Casimir-gap}.  Thus \(C=0\) and
\(\xi\in W\rho\).
\end{proof}

Thus \(J(\mu)\) lies in the regular minimal primitive fibre.
Petukhov proves that for a simply-laced simple Lie algebra not of type \(A\),
this fibre contains
exactly \(\ell\) primitive ideals, indexed by Dynkin vertices
\cite[Proposition~5.4]{Petukhov2018}.  Put
\[
   J_i:=\operatorname{Ann}_{U(\mathfrak g)}L_{\mathfrak g}(-\alpha_i).
\]
With Petukhov's convention
\(\tau_R(w)=\{\alpha\in\Pi:w\alpha\in\Phi^+\}\), these are precisely the
ideals with \(\tau(J_i)=\Pi\setminus\{\alpha_i\}\).

To compute their finite-\(W\) highest weights, set \(y_i=z_i^{-1}\) and
\(\nu_i=y_i\rho-\rho\).  The path word \(y_i\) begins with \(s_2\) and ends
with \(s_i\), so its ordinary Coxeter left and right descent sets are
\(\{s_2\}\) and \(\{s_i\}\).  Hence
\[
   \tau\!\left(\operatorname{Ann}L_{\mathfrak g}(\nu_i)\right)
   =\Pi\setminus\{\alpha_i\},
   \qquad
   \operatorname{Ann}L_{\mathfrak g}(\nu_i)=J_i,
\]
by \cite[Proposition~5.4]{Petukhov2018}; this is the reason for the inverse.
Moreover \(y_i^{-1}\alpha_2<0\), so \(\nu_i\) is the
\(\langle s_2\rangle\)-antidominant representative for the
\(\alpha_2\)-Whittaker model.  Lemma~\ref{lem:primitive-finiteW-bridge}
therefore applies with this prescribed \(\nu_i\) and gives
\[
   \lambda_i^H=\bar\delta+\overline{\nu_i},
   \qquad C=(\nu_i,\nu_i+2\rho)=0.
\]
The path calculation is recorded in Appendix~\ref{app:type-D-calculations}.

\begin{proposition}[Regular node weights]
The finite-\(W\) highest weights attached to \(J_i\) are
\begin{align}
   \lambda_1^H
   &=
   r\varpi,
   \label{eq:node-weight-1}\\
   \lambda_2^H
   &=
   (r-1)\varpi+\eta_1,
   \\
   \lambda_i^H
   &=
   (r+1-i)\varpi+\eta_{i-1},
   &&3\le i\le r-1,
   \label{eq:node-weight-middle}\\
   \lambda_{\ell-2}^H
   &=
   \varpi+\eta_{r-1}+\eta_r,
   \\
   \lambda_{\ell-1}^H
   &=
   2\eta_r,
   \\
   \lambda_\ell^H
   &=
   2\eta_{r-1}.
   \label{eq:node-weight-spin2}
\end{align}
For \(r=3\), the middle range is empty and the last three formulas are read
under \(D_3\cong A_3\).
\end{proposition}

\begin{proposition}[Node-\(2\) rigidity]
\label{prop:node2-rigidity}
If
\[
   I_\ell^{\mathrm{cand}}\subset J_i
\]
for a regular minimal-orbit primitive ideal \(J_i\), then \(i=2\).
\end{proposition}

\begin{proof}
If \(I_\ell^{\mathrm{cand}}\subset J_i\), then the corresponding
finite-dimensional simple \(H\)-module is a \(B_\ell\)-module.

For \(i=1\), the \(A_1\)-highest weight in
\eqref{eq:node-weight-1} is \(r\varpi\).  This contradicts the
\(B_\ell\)-bound
\[
   a\le r-1
\]
from Lemma~\ref{lem:allowed-finite-types}.

For \(i=2\),
\[
   \lambda_2^H
   =
   (r-1)\varpi+\eta_1
\]
satisfies both bounds.

For every \(i\ge3\), the \(D_r\)-component in
\eqref{eq:node-weight-middle}--\eqref{eq:node-weight-spin2} satisfies
\[
   \left\langle
      (\lambda_i^H)_D,\theta_D^\vee
   \right\rangle
   =
   2.
\]
This contradicts \(e_D^2=0\), equivalently
\[
   \left\langle
      \lambda_D,\theta_D^\vee
   \right\rangle
   \le1.
\]
The same computation holds for \(r=3\cong A_3\); it is spelled out in
Appendix~\ref{app:rank-five-boundary}.  Hence only \(i=2\) survives.
\end{proof}

\begin{proposition}[Minimal-orbit rigidity]
\label{prop:minimal-orbit-rigidity}
Under \eqref{eq:minimal-orbit-assumptions},
\[
   J(\mu)=J_{2,\ell},
\]
and
\[
   \mu\in\{\mu_1,\ldots,\mu_\ell\}.
\]
\end{proposition}

\begin{proof}
Proposition~\ref{prop:minimal-regular-central-character} places \(J(\mu)\)
in the regular integral minimal-orbit fibre.  Hence
\[
   J(\mu)=J_i
\]
for some \(i\).  Proposition~\ref{prop:node2-rigidity} forces \(i=2\).

It remains to determine the highest weights with annihilator \(J_{2,\ell}\).
At regular integral central character, equality of highest-weight primitive
ideals is the Kazhdan--Lusztig left-cell relation.  The finite subregular
left cell containing \(s_2\) consists exactly of the unique reduced path
words ending at \(s_2\):
\[
   z_1,z_2,\ldots,z_\ell.
\]
See \cite[Chapter~12]{Bonnafe2017} and
\cite[Proposition~5.4]{Petukhov2018}.  Therefore
\[
   \mu+\rho=z_i\rho
\]
for a unique \(1\le i\le\ell\), and hence
\[
   \mu=z_i\rho-\rho=\mu_i.
\]
\end{proof}

\subsection{Completion of exhaustion}

The two associated-variety cases give the exhaustion statement.

\begin{theorem}[Exhaustion theorem]
\label{thm:exhaustion}
Let \(\ell\ge5\).  If
\[
   I_\ell^{\mathrm{cand}}
   \subset
   \operatorname{Ann}_{U(D_\ell)}L_{D_\ell}(\mu)
\]
and
\[
   \mu+\rho
   \in
   W(D_\ell)\rho
   +(2\ell-3)Q(D_\ell),
\]
then
\begin{equation}
   \mu\in\{\mu_0,\mu_1,\ldots,\mu_\ell\}.
\end{equation}
More precisely,
\[
   \mu=0
\]
in the zero-orbit branch, while
\[
   \mu\in\{\mu_1,\ldots,\mu_\ell\}
\]
in the minimal-orbit branch.
\end{theorem}

\begin{proof}
Lemma~\ref{lem:filtered-Zhu-containment} gives
\[
   \mathcal V(J(\mu))
   =
   \{0\}
   \quad\text{or}\quad
   \overline{\mathcal O_{\min}}.
\]
In the first case,
Proposition~\ref{prop:zero-orbit-rigidity} gives \(\mu=0=\mu_0\).
In the second case,
Proposition~\ref{prop:minimal-orbit-rigidity} gives
\(\mu\in\{\mu_1,\ldots,\mu_\ell\}\).
\end{proof}

\section{The affine left-cell conjecture at level -1}
\label{sec:final-left-cell}

By Propositions~\ref{prop:subregular-left-cell} and
\ref{prop:path-dot-action},
\[
   \mathbf c^L(s_0)=\{w_0,w_1,\ldots,w_\ell\},
   \qquad
   w_i\circ(-\Lambda_0)=-\Lambda_0+\mu_i.
\]
Here \(m=k+h^\vee=N=h^\vee-1\).  In type~\(D_\ell\),
\cite[Section~2.4.7]{ShanYanZhao2025} gives
\(\check{\mathbb O}(m)=\check{\mathbb O}_{\mathrm{sreg}}\), a distinguished
orbit.  Since \(A_\ell\) is dominant with stabilizer \(\langle s_0\rangle\),
we have \(\xi_m=A_\ell\) and \(w_m=s_0\).  Thus
\cite[Conjecture~5.2.1]{ShanYanZhao2026} specializes to the assertion proved here.

\subsection{Classification and passage to the simple quotient}

\begin{theorem}[Candidate quotient classification]
\label{thm:classification-Q}
The simple objects in the affine vacuum block that factor through \(Q_\ell\)
are exactly
\begin{equation}
   \left\{
      L(-\Lambda_0+\mu_i)\;\middle|\;0\le i\le\ell
   \right\}.
\end{equation}
\end{theorem}

\begin{proof}
The vacuum is present.  For \(i\ge1\), Theorems
\ref{thm:membership} and \ref{thm:primitive-annihilator-compression} give
\[
   I_\ell^{\mathrm{cand}}
   \subset J_{2,\ell}
   =\operatorname{Ann}_{U(D_\ell)}L_{D_\ell}(\mu_i),
\]
so all path modules descend.  Conversely, by the classical affine
linkage/block parametrization \cite{Kac1990}
\cite[Section~1.2]{ShanYanZhao2026}, any simple object in the ambient block of
\(-\Lambda_0\) is \(L(y\circ(-\Lambda_0))\) for some \(y\in\widehat W\),
modulo the usual stabilizer redundancy.

If such a simple factors through \(Q_\ell\), its degree-zero top is
annihilated by \(I_\ell^{\mathrm{cand}}\).  Write its finite highest-weight
part as \(\mu\).  Since \(\widehat W=W\ltimes Q^\vee\) and a
translation by \(q\) adds \(Nq\) to the finite part of a level-\(N\)
weight \cite[Chapter~6]{Kac1990}, the shifted vacuum
\(-\Lambda_0+\widehat\rho=N\Lambda_0+\rho\) gives the necessary
congruence
\(\mu+\rho\in W\rho+(2\ell-3)Q(D_\ell)\), because type \(D\) is
simply laced and \(Q^\vee=Q\).  Theorem~\ref{thm:exhaustion} therefore
forces \(\mu\in\{\mu_0,\ldots,\mu_\ell\}\).
\end{proof}

The type-\(D_\ell\), level-\(-1\) maximal-ideal theorem of
\cite{Jin2026} states
\begin{equation}
\label{eq:Q-is-simple-final}
   Q_\ell\cong L_{-1}(D_\ell).
\end{equation}
Simplicity of the candidate quotient is used only here.

\begin{theorem}[Simple objects]
\label{thm:final-simple-classification}
For every \(\ell\ge5\),
\begin{equation}
   \operatorname{Irr}\mathcal O_{-\Lambda_0}
   \!\left(L_{-1}(D_\ell)\right)
   =
   \left\{
      L\!\left(w_i\circ(-\Lambda_0)\right)
      \;\middle|\;0\le i\le\ell
   \right\}.
\end{equation}
In particular, the block has exactly \(\ell+1\) simple objects.
\end{theorem}

\begin{proof}
Combine Theorem~\ref{thm:classification-Q}, \eqref{eq:Q-is-simple-final},
and Proposition~\ref{prop:path-dot-action}.
\end{proof}

\subsection{Grothendieck group}

The ambient realization in \cite[Section~5.2]{ShanYanZhao2026} embeds
\(K_0\mathcal O_{-\Lambda_0}(L_{-1}(D_\ell))\) into the specialized affine
Hecke quotient, sending \([L(y\circ(-\Lambda_0))]\) to \(D_y|_{q=1}\).
On the Hecke side, the canonical inclusion
\cite[Equation~(2.1.2)]{ShanYanZhao2026} embeds the dual left-cell module in
the same quotient, with image spanned by the \(D_y|_{q=1}\) for
\(y\in\mathbf c^L(s_0)\).  Since the simple labels are exactly the elements
of this left cell, the two canonical subspaces agree.

\begin{theorem}[Shan--Yan--Zhao conjecture]
For every \(\ell\ge5\), under the canonical ambient embedding above,
\begin{equation}
\label{eq:final-SYZ-K0}
   K_0\mathcal O_{-\Lambda_0}
   \!\left(L_{-1}(D_\ell)\right)
   \cong
   \left.
      \mathcal H^\vee_{\mathrm{aff},\mathbf c^L(s_0)}
   \right|_{q=1}.
\end{equation}
Hence \cite[Conjecture~5.2.1]{ShanYanZhao2026} holds for
\((D_\ell,-1)\).
\end{theorem}

\begin{proof}
By Theorem~\ref{thm:final-simple-classification} and
Proposition~\ref{prop:subregular-left-cell}, the image of the left-hand side
in the ambient quotient is exactly
\(\operatorname{span}_{\mathbb Z}\{D_y|_{q=1}:y\in\mathbf c^L(s_0)\}\).
By the canonical cell inclusion just recalled, this is precisely the image
of the specialized dual left-cell module, proving
\eqref{eq:final-SYZ-K0}.  The image is therefore
\(\widehat W\)-stable, so the ambient action restricts to it without requiring
twisting functors to preserve the vertex-algebra subcategory.
\end{proof}

\subsection{Closed character formulas}

Let \(\mathbf K\) be the affine central element,
\(a_i=\langle\Lambda_i,\mathbf K\rangle\), so
\[
   a_0=a_1=a_{\ell-1}=a_\ell=1,
   \qquad
   a_2=\cdots=a_{\ell-2}=2.
\]
Let \(R_{\widehat{\mathfrak g}}\) be the affine Weyl denominator, let
\(t_\gamma\) denote translation by \(\gamma\in Q(D_\ell)\), put
\(A_\ell=N\Lambda_0+\rho\), and set
\[
   b_i(u,\gamma)
   :=\langle\Lambda_i,u\theta^\vee\rangle
      -a_i\bigl((\theta,\gamma)+1\bigr).
\]

\begin{corollary}[Uniform characters]
\label{cor:uniform-characters}
For \(0\le i\le\ell\),
\begin{equation}
\label{eq:uniform-character}
   R_{\widehat{\mathfrak g}}\,
   \operatorname{ch}L\!\left(w_i\circ(-\Lambda_0)\right)
   =
   \frac{\varepsilon(w_i)}2
   \sum_{u\in W(D_\ell)}\varepsilon(u)
   \sum_{\gamma\in Q(D_\ell)}
   b_i(u,\gamma)e^{u t_\gamma A_\ell}.
\end{equation}
\end{corollary}

\begin{proof}
By \cite[Proposition~3.6 and Lemma~A.3]{BezrukavnikovKacKrylov2024},
\(R_{\widehat{\mathfrak g}}\operatorname{ch}L(w_i\circ(-\Lambda_0))
=\sum_x\varepsilon(xw_i)m^x_{w_i}e^{xA_\ell}\).  Since
\(s_0A_\ell=A_\ell\), pair \(x=ut_\gamma\) with \(xs_0\).  As
\(s_0=t_\theta s_\theta\), their right \(W\)-cosets are
\(t_{u\gamma}W\) and \(t_{u(\gamma+\theta)}W\); hence
\cite[Proposition~A.6 and Equation~(31)]{BezrukavnikovKacKrylov2024} yields
\(m^x_{w_i}-m^{xs_0}_{w_i}=b_i(u,\gamma)\), using
\(Q^\vee=Q\) and \(\lVert\theta\rVert^2=2\).  Pairing contributes
\(1/2\), and \(\varepsilon(t_\gamma)=1\), proving
\eqref{eq:uniform-character}.
\end{proof}

\appendix
\section{Type-D root calculations}
\label{app:type-D-calculations}

This appendix records the coordinate calculations used in
Sections~\ref{sec:membership} and \ref{sec:exhaustion}.  Put
\(r=\ell-2\), use orthogonal coordinates
\(\varepsilon_1,\ldots,\varepsilon_\ell\), and work in Premet's minimal
root model
\[
   \beta=\alpha_2=\varepsilon_2-\varepsilon_3.
\]
The roots orthogonal to \(\beta\) form \(A_1\sqcup D_r\).  Take
\[
   \phi=\varepsilon_2+\varepsilon_3,
   \qquad
   \varpi_\beta=\frac{\phi}{2},
\]
and for the \(D_r\)-factor use
\begin{equation}
   \delta_1=\varepsilon_1,
   \qquad
   \delta_j=\varepsilon_{j+2}\quad(2\le j\le r).
\end{equation}
For \(r\ge4\), its fundamental weights are
\[
   \eta_j=\delta_1+\cdots+\delta_j\quad(1\le j\le r-2),
\]
\[
   \eta_{r-1}=\tfrac12(\delta_1+\cdots+\delta_{r-1}-\delta_r),
   \qquad
   \eta_r=\tfrac12(\delta_1+\cdots+\delta_{r-1}+\delta_r).
\]
We transport \(\varpi_\beta\) to the \(A_1\)-weight \(\varpi\) in the
highest-root model.

\subsection{Premet's shift}

Premet's shift is the restriction of
\(\delta=\frac12\sum_{\gamma\in\Phi^+_{-1}}\gamma\).  The positive roots
with \(\langle\gamma,\beta^\vee\rangle=-1\) are
\[
   \varepsilon_1-\varepsilon_2,
   \quad
   \varepsilon_1+\varepsilon_3,
   \quad
   \varepsilon_3\pm\varepsilon_j\quad(4\le j\le\ell).
\]
Set
\[
 x_*^\beta=\varepsilon_1+\frac{\varepsilon_2+\varepsilon_3}{2}.
\]
The positive \(A_1\oplus D_r\)-roots have \(x_*^\beta\)-charges in
\(\{0,1\}\), while the three displayed root families have charges
\(\frac12,\frac32,\frac12\), respectively.  Since
\(\beta(x_*^\beta)=0\), the corresponding positive roots
\(\beta+\gamma\in\Phi^+_{e,1}\) have the same charges.
Their sum is
\[
   2\varepsilon_1-\varepsilon_2+(2\ell-5)\varepsilon_3,
\]
so
\[
   \delta=
   \varepsilon_1-\frac12\varepsilon_2
   +\frac{2\ell-5}{2}\varepsilon_3.
\]
Pairing with \(\phi^\vee\) gives \(\ell-3=r-1\), while the
\(D_r\)-restriction is \(\delta_1=\eta_1\).  Therefore
\begin{proposition}[Premet shift in type \(D\)]
\begin{equation}
\label{eq:appendix-D-Premet-shift}
   \bar\delta=(r-1)\varpi+\eta_1.
\end{equation}
\end{proposition}
This proves Lemma~\ref{lem:Premet-shift-D}.

\subsection{Inverse path words and regular node weights}

Let \(y_i=z_i^{-1}\), as fixed in Section~\ref{sec:exhaustion}; the
inverse is the one that has right descent \(s_i\) and hence labels the node
ideal \(J_i\).  Applying it to
\(\rho=(\ell-1,\ell-2,\ldots,1,0)\) gives, with \(\nu_i=y_i\rho-\rho\),
\begin{align}
 \nu_1&=-\varepsilon_1-\varepsilon_2+2\varepsilon_3,
 &\nu_2&=-\varepsilon_2+\varepsilon_3,\label{eq:appendix-D-nu12}\\
 \nu_i&=-(i-1)\varepsilon_2+
        \varepsilon_3+\cdots+\varepsilon_{i+1}
 &&(3\le i\le\ell-2),\\
 \nu_{\ell-1}&=-r\varepsilon_2+
       \varepsilon_3+\cdots+\varepsilon_\ell,\\
 \nu_\ell&=-r\varepsilon_2+
       \varepsilon_3+\cdots+\varepsilon_{\ell-1}-\varepsilon_\ell.
       \label{eq:appendix-D-nusp2}
\end{align}
Premet's comparison theorem gives
\[
   \lambda_i^H=\bar\delta+\overline{\nu_i}.
\]
For an ambient weight \(\sum c_j\varepsilon_j\), its \(A_1\)-coefficient
in units of \(\varpi\) is \(c_2+c_3\), while its \(D_r\)-restriction keeps
coordinates \(c_1,c_4,\ldots,c_\ell\).  Substitution yields the following.

\begin{proposition}[Node-weight calculation]
\begin{align}
 \lambda_1^H&=r\varpi,
 &\lambda_2^H&=(r-1)\varpi+\eta_1,\\
 \lambda_i^H&=(r+1-i)\varpi+\eta_{i-1}
 &&(3\le i\le r-1),\\
 \lambda_{\ell-2}^H&=\varpi+\eta_{r-1}+\eta_r,\\
 \lambda_{\ell-1}^H&=2\eta_r,
 &\lambda_\ell^H&=2\eta_{r-1}.
\end{align}
For the highest root \(\theta_D=\delta_1+\delta_2\),
\[
   \langle(\lambda_2^H)_D,\theta_D^\vee\rangle=1,
   \qquad
   \langle(\lambda_i^H)_D,\theta_D^\vee\rangle=2
   \quad(3\le i\le\ell).
\]
\end{proposition}

\begin{proof}
Adding \eqref{eq:appendix-D-Premet-shift} to
\eqref{eq:appendix-D-nu12}--\eqref{eq:appendix-D-nusp2} gives the displayed
weights.  The last assertion follows by pairing the \(D_r\)-parts with
\(\delta_1+\delta_2\); at the branch and spinor endpoints use
\[
   \eta_{r-1}+\eta_r=\delta_1+\cdots+\delta_{r-1},
   \qquad
   2\eta_{r\,\text{ or }\,r-1}
   =\delta_1+\cdots+\delta_{r-1}\pm\delta_r.
\]
\end{proof}

For \(r\ge4\), the root lattice is
\[
   Q(D_r)=\left\{\sum n_j\delta_j\;\middle|\;
   n_j\in\mathbb Z,\ \sum n_j\in2\mathbb Z\right\},
\]
and
\[
   Q(D_r)+\mathbb Z\eta_1
   =Q(D_r)\sqcup(\eta_1+Q(D_r)).
\]
This root/vector subgroup is used in the Casimir estimates.  The
boundary case \(r=3\cong A_3\) is recorded next.
\section{The rank-five boundary}
\label{app:rank-five-boundary}

When \(\ell=5\), one has \(r=3\), \(N=7\), and \(D_3\cong A_3\).
There is no low-rank degeneration in the candidate relations: the second
Kostant component has highest weight \(\omega_4+\omega_5\), and
\[
   [s_D]_{\mathrm{DS}}=c_D:J_D^2:\ne0,
\]
so the relation \(e_D^2=0\) used in Lemma~\ref{lem:allowed-finite-types}
remains valid \cite[Equation~(32), Proposition~3.6, and Remark~3.15]{Jin2026}.

In orthogonal coordinates \(\delta_1,\delta_2,\delta_3\), take
\[
\begin{gathered}
 \beta_1=\delta_1-\delta_2,\quad
 \beta_2=\delta_2-\delta_3,\quad
 \beta_3=\delta_2+\delta_3,\\
 \eta_1=\delta_1,\quad
 \eta_2=\tfrac12(\delta_1+\delta_2-\delta_3),\quad
 \eta_3=\tfrac12(\delta_1+\delta_2+\delta_3).
\end{gathered}
\]
Under \(D_3\cong A_3\), these are respectively
\(\omega_2,\omega_1,\omega_3\).  Since
\(\theta_D=\delta_1+\delta_2\), all three satisfy
\(\langle\eta_j,\theta_D^\vee\rangle=1\).  Moreover
\[
\begin{aligned}
 Q(D_3)&=\{n\in\mathbb Z^3\mid n_1+n_2+n_3\in2\mathbb Z\},\\
 Q(D_3)+\mathbb Z\eta_1&=Q(D_3)\sqcup(\eta_1+Q(D_3)).
\end{aligned}
\]
The second line is the order-two subgroup generated by \(\omega_2\) in
\(P(A_3)/Q(A_3)\cong\mathbb Z/4\mathbb Z\).  Hence the level-one
root/vector-coset argument leaves only \(0\) and \(\eta_1\); this covers
Lemmas~\ref{lem:zero-root-vector-bound} and \ref{lem:minimal-no-spinor}.

With \(\rho_{D_3}=2\delta_1+\delta_2\), one has
\(c_2(\eta_1)=5\), \(C_{\max}=49/6\), and \(C\le4/3<2\), so
Proposition~\ref{prop:minimal-Casimir-gap} holds.  The node weights are
\(3\varpi,2\varpi+\eta_1,\varpi+\eta_2+\eta_3,2\eta_3,2\eta_2\): node~1
violates \(a\le2\), whereas nodes~3,4,5 have highest-root level~2, proving
Proposition~\ref{prop:node2-rigidity}.  Finally
\(s^2\le12a-a^2+50/3-10c\le110/3<49=N^2\), so the zero-orbit argument is
unchanged.  No exceptional modification is therefore required for \(D_5\).

\providecommand{\bysame}{\leavevmode\hbox to3em{\hrulefill}\thinspace}
\providecommand{\MR}{\relax\ifhmode\unskip\space\fi MR }
% \MRhref is called by the amsart/book/proc definition of \MR.
\providecommand{\MRhref}[2]{%
  \href{http://www.ams.org/mathscinet-getitem?mr=#1}{#2}
}
\providecommand{\href}[2]{#2}
\begin{thebibliography}{KMFP25}

\bibitem[AM18]{ArakawaMoreau2018}
Tomoyuki Arakawa and Anne Moreau, \emph{Joseph ideals and lisse minimal
  {$W$}-algebras}, Journal of the Institute of Mathematics of Jussieu
  \textbf{17} (2018), no.~2, 397--417.

\bibitem[Ara05]{Arakawa2005}
Tomoyuki Arakawa, \emph{Representation theory of superconformal algebras and
  the {Kac--Roan--Wakimoto} conjecture}, Duke Mathematical Journal \textbf{130}
  (2005), no.~3, 435--478.

\bibitem[Ara15]{Arakawa2015}
\bysame, \emph{Rationality of {$W$}-algebras: principal nilpotent cases},
  Annals of Mathematics \textbf{182} (2015), no.~2, 565--604.

\bibitem[BKK24]{BezrukavnikovKacKrylov2024}
Roman Bezrukavnikov, Victor~G. Kac, and Vasily Krylov, \emph{Subregular
  nilpotent orbits and explicit character formulas for modules over affine lie
  algebras}, Pure and Applied Mathematics Quarterly \textbf{20} (2024), no.~1,
  81--138.

\bibitem[Bon17]{Bonnafe2017}
C{\'e}dric Bonnaf{\'e}, \emph{Kazhdan--lusztig cells with unequal parameters},
  Algebra and Applications, vol.~24, Springer, Cham, 2017.

\bibitem[Che21]{Chen2021Annihilator}
Chih-Whi Chen, \emph{Annihilator ideals and blocks of whittaker modules over
  quasireductive lie superalgebras}, 2021.

\bibitem[DLM98]{DongLiMason1998}
Chongying Dong, Haisheng Li, and Geoffrey Mason, \emph{Vertex operator algebras
  and associative algebras}, Journal of Algebra \textbf{206} (1998), no.~1,
  67--96.

\bibitem[Jin26]{Jin2026}
Sihai Jin, \emph{Completing the {Arakawa--Moreau} conjecture on maximal ideals
  of affine vertex algebras}, arXiv:2607.25249, 2026.

\bibitem[Jos95]{Joseph1995}
Anthony Joseph, \emph{Quantum groups and their primitive ideals}, Ergebnisse
  der Mathematik und ihrer Grenzgebiete, 3. Folge, vol.~29, Springer-Verlag,
  Berlin, 1995.

\bibitem[Kac90]{Kac1990}
Victor~G. Kac, \emph{Infinite-dimensional lie algebras}, 3 ed., Cambridge
  University Press, Cambridge, 1990.

\bibitem[KMFP25]{KacMosenederFrajriaPapi2025}
Victor~G. Kac, Pierluigi M\"oseneder~Frajria, and Paolo Papi, \emph{Unitarity
  of minimal {$W$}-algebras and their representations {II}: {Ramond} sector},
  Japanese Journal of Mathematics \textbf{20} (2025), no.~2, 169--281.

\bibitem[KW04]{KacWakimoto2004}
Victor~G. Kac and Minoru Wakimoto, \emph{Quantum reduction and representation
  theory of superconformal algebras}, Advances in Mathematics \textbf{185}
  (2004), no.~2, 400--458.

\bibitem[Li96]{Li1996}
Haisheng Li, \emph{Local systems of twisted vertex operators, vertex operator
  superalgebras and twisted modules}, Moonshine, the Monster, and Related
  Topics, Contemporary Mathematics, vol. 193, American Mathematical Society,
  Providence, RI, 1996, pp.~203--236.

\bibitem[MS97]{MilicicSoergel1997}
Dragan Mili\v{c}i\'c and Wolfgang Soergel, \emph{The composition series of
  modules induced from whittaker modules}, Commentarii Mathematici Helvetici
  \textbf{72} (1997), no.~4, 503--520.

\bibitem[Pet18]{Petukhov2018}
Alexey Petukhov, \emph{Finite-dimensional representations of minimal nilpotent
  {$W$}-algebras and zigzag algebras}, Representation Theory \textbf{22}
  (2018), 223--245.

\bibitem[Pre07]{Premet2007}
Alexander Premet, \emph{Enveloping algebras of slodowy slices and the joseph
  ideal}, Journal of the European Mathematical Society \textbf{9} (2007),
  no.~3, 487--543.

\bibitem[SYZ25]{ShanYanZhao2025}
Peng Shan, Wenbin Yan, and Qixian Zhao, \emph{Cyclotomic level maps and
  associated varieties of simple affine vertex algebras}, arXiv:2507.09254, 2025.

\bibitem[SYZ26]{ShanYanZhao2026}
\bysame, \emph{Affine vertex algebras and an affine analog of
  {Barbasch--Vogan}'s construction}, arXiv:2608.00428, 2026.

\bibitem[Xu19]{Xu2019}
Tianyuan Xu, \emph{On the subregular {$J$}-rings of coxeter systems}, Algebras
  and Representation Theory \textbf{22} (2019), no.~6, 1479--1512.

\end{thebibliography}

\end{document}
